% SGA 3, 第一卷；原 PDF 第 1--26 页，第 I 讲第 4 节标题之前。
\SourcePage{1}
\SGAsection{}{再版说明}
这里呈献的是原讨论班讲义略加修订的再版；其宗旨和内容见引言。修订主要包括：订正排印错误，以脚注增补若干注记或参考文献，将全书划为目前的三卷，每卷配有详细目录和记号索引，并在第三卷末附上术语索引。此外，J.-E. Bertin 亲自重写了第 VI\textsubscript{B} 讲的一部分，尤其是第 5 节和第 10 节；因此，有些对该讲的引用与原版中的引用不同。本卷卷首列有讨论班各讲的清单。

本讨论班第一版问世以后，《\Term{algebraic-geometry}{代数几何}基础》第 IV 章已全部出版，使本讲义的某些段落成为多余；我们有时在脚注中指出 EGA IV 的有关出处，读者可借此略过这些段落。

关于另一部系统采用\Term{scheme}{概形}语言讲解\Term{algebraic-group}{代数群}的著作，我们推荐 M. Demazure 与 P. Gabriel 合著的 \emph{Groupes Algébriques}（North-Holland 与 Masson et Cie）。这部书不预设任何\Term{algebraic-geometry}{代数几何}知识，而包含全部必需的\Term{scheme}{概形}论预备知识，因此可作为学习本讨论班讲义的入门书。（它还讨论本讲义未涉及的主题，例如该书第 V 章中的交换\Term{affine-algebraic-group}{仿射代数群}的 Dieudonné 式结构理论。）

\begin{flushright}Bures-sur-Yvette，1970 年 3 月\end{flushright}
\SourcePage{2}
% 原 PDF 第 2 页为空白页。
\SourcePage{3}
\SGAsection{}{引言}
\textbf{1.} 本讨论班有双重目的。

一方面，我们希望为一般的\Term{group-scheme}{群概形}理论提供便利的基础。第 I 至 IV 讲将给出为此不可缺少的\Term{scheme}{概形}及\Term{category}{范畴}语法的初步练习。为了得到一种能够毫不费力地贴合几何直观的语言，并避免长期使用起来令人难以忍受的迂回说法，我们始终把另一个\Term{prescheme}{预概形} $S$ 上的\Term{prescheme}{预概形} $X$ 与它所\Term{representable-functor}{代表的函子} $(\mathrm{Sch})_{/S}^{\circ}\to(\mathrm{Ens})$ 等同起来\footnote{这一观点似乎是 P. Cartier 在八九年前研究\Term{formal-group}{形式群}理论时首先想到的；遗憾的是，他没有花力气将它作出应有的精确表述和系统整理。}；因此，许多定义必须表述为适用于任意\Term{functor}{函子}，无论其\Term{representable}{可表}与否。事实上，我们要用的几乎所有\Term{functor}{函子}都是“\Term{sheaf}{层}”（对于“\Term{faithfully-flat}{忠实平坦}\Term{quasi-compact}{拟紧}\Term{topology}{拓扑}”而言）；第 IV 讲只附带涉及群，勾勒了“\Term{localization}{局部化}”和层的语言，这种语言在\Term{functor}{函子}的\Term{representability}{可表性}问题中也极为便利。尤其对于取商问题，该讲将为后文提供最方便的框架。

第 V 讲给出关于商的存在性的一些一般结果；第 VI\textsubscript{A} 讲又在域上（或更一般地，在 Artin 环上）的\Term{algebraic-group}{代数群}对其\Term{subgroup}{子群}取商的情形中重述这些结果\footnote{关于取商的更深入研究，尤其关于\Term{reductive-group}{约化群}的商，见 D. Mumford 的重要著作 \emph{Geometric Invariant Theory}，Ergebnisse der Mathematik，第 34 卷，Springer，1965。须指出，该书的术语在一个重要方面与我们的不同：在\Term{characteristic}{特征} $p>0$ 的域上，Mumford 称为“\Term{reduced}{约化}”的群\textsuperscript{(1)}，正是本讨论班意义下\Term{smooth}{光滑}的\Term{multiplicative-type}{乘法型群}（参见第 IX 讲）。Mumford 对“\Term{reduced}{约化}”一词的用法在底不为域时失去了意义；可以认为，他只是暂时采用这种用法，将它当作一种权宜之计（他在序言第二段出于别的理由所作的说明，也大致如此！）。\textsuperscript{(1)} 编者注：见本引言末增补的注记。}。这后一讲和紧随其后的第 VI\textsubscript{B} 讲，还包含只对域上\Term{algebraic-group}{代数群}成立、并将在后文经常使用的若干初等结果。

第 VII 讲研究与\Term{base-field}{基域}\Term{characteristic}{特征}有关的若干事实，尤其以适当的一般性阐述\Term{height}{高度}为 $1$ 的\Term{radicial}{根性}\Term{group-scheme}{群概形}与\Term{restriction}{限制} $p$-\Term{lie-algebra}{Lie 代数}之间的对应。
\SourcePage{4}
最后，第 XVIII 讲把 Weil 关于\Term{algebraic-group}{代数群}的“\Term{birational}{双有理}”定义的\Term{theorem}{定理}推广到\Term{scheme}{概形}论。

另一方面，我们打算将\Term{affine-algebraic-group}{仿射代数群}的 Borel--Chevalley 结构理论推广到任意底\Term{prescheme}{预概形}上的群。在撰写讨论班笔记时还发现，\Term{affine}{仿射}性假设对于该理论的许多结果并不必要。最完整的结果自然是在“\Term{semisimple}{半单}\Term{group-scheme}{群概形}”，或更一般的“\Term{reductive}{约化}\Term{group-scheme}{群概形}”的情形中得到的；从第 XIX 讲起，我们将专门研究它们。Chevalley 本人已在整数环上构造了“Tohoku 群”，本讨论班将重新给出这一构造。主要的唯一性\Term{theorem}{定理}\textsuperscript{(2)} 给出了这些 Tohoku 群在底\Term{prescheme}{预概形} $S$ 上的“\Term{twisting}{扭曲}”变体的一个简单刻画：它们是 $S$ 上的\Term{affine}{仿射}\Term{smooth}{光滑}群，其\Term{geometric-fibre}{几何纤维}是通常意义下的\Term{connected}{连通}\Term{semisimple-group}{半单群}\footnote{这正是 M. Demazure 博士论文的主要结果（\emph{Schémas en groupes réductifs}，Bull. Soc. Math. France \textbf{93}（1965），369--413）。}\footnote{\textsuperscript{(2)} 编者注：这里指推论 XXIII.5.6，它容易由\Term{split-reductive-group}{分裂约化群}的唯一性\Term{theorem}{定理}\Term{pushout}{推出}（参见 XXIII，\Term{theorem}{定理} 4.1 和推论 5.2），因为每个\Term{semisimple}{半单} $S$-群都是一个“Tohoku 群”的（对于\Term{etale-topology}{平展拓扑}的）“\Term{twisting}{扭曲}\Term{form}{形式}”（参见 XXII，推论 2.3）。}。

\textbf{2.} 与\Term{algebraically-closed-field}{代数闭域}上的已知理论一样，所考察的\Term{group-scheme}{群概形}的\Term{subtorus}{子环面}起着关键\Term{action}{作用}。因此，\Term{torus}{环面}以及更一般的“\Term{multiplicative-type}{乘法型}\Term{group-scheme}{群概形}”的预备研究，在本讨论班中占据相当大的篇幅（第 VIII 至 XII 讲），既从内在的观点，也从给定群的\Term{multiplicative-type}{乘法型}\Term{subgroup}{子群}的观点加以研究。它们显著的\Term{rigidity}{刚性}（在某些方面甚至比\Term{abelian-scheme}{阿贝尔概形}或\Term{semisimple-group-scheme}{半单群概形}的\Term{rigidity}{刚性}还强）使它们成为研究某些更一般群的十分\Term{effective}{有效的}工具。

\textbf{3.} 从第 XII 讲开始（已提到的第 XVIII 讲除外），我们将经常使用\Term{algebraically-closed-field}{代数闭域}上的\Term{affine-algebraic-group}{仿射代数群}理论；读者可在 1956 年 Chevalley 讨论班讲义中，尤其在其第 IV 至 IX 讲中找到这一理论。我们也将使用该讨论班后续各讲关于\Term{semisimple}{半单}\Term{algebraic-group}{代数群}结构的内容，不过程度较轻。事实上，我们将直接在\Term{scheme}{概形}框架内重新建立 Chevalley 理论：我们将看到，这种做法（即使在底为域时）也使论述更加简单而精确。

\textbf{4.} 本讨论班的主要目的是发展适用于任意底上\Term{group-scheme}{群概形}的技术，换言之，本质上是研究\Term{algebraic-group}{代数群}族的技术。因此，这些族的\Term{infinitesimal}{无穷小}性质，尤其是\Term{base-scheme}{底概形}为 Artin \Term{scheme}{概形}的情形，起着重要\Term{action}{作用}。即使研究域 $k$ 上的\Term{algebraic-group}{代数群}，当该域不是\Term{perfect-field}{完美域}时，
\SourcePage{5}
也要用到这些性质，特别是为了在非 Galois 情形应用\Term{descent}{下降}技术。在这种情形得到的新结果中，我们指出：任意\Term{smooth}{光滑}\Term{algebraic-group}{代数群}中\Term{maximal-torus}{极大环面}和 \Term{cartan-subgroup}{Cartan 子群}的存在性、\Term{maximal-torus}{极大环面}簇的有理性及各种相关结果（第 XIV 讲\footnote{\textsuperscript{(3)} 编者注：尤其是\Term{theorem}{定理} 1.1 和 6.1。}），以及\Term{semisimple-group}{半单群}的“\Term{form}{形式}”与某个适当的\Term{semisimple}{半单}\Term{algebraic-group}{代数群}（一般不\Term{connected}{连通}）下的\Term{principal-homogeneous-bundle}{主齐性纤维丛}之间的对应（第 XXIV 讲，1.16--1.20）。一般而言，在非\Term{perfect}{完美}\Term{base-field}{基域}上工作所需的方法，本质上就是处理任意底\Term{prescheme}{预概形}所用的方法，因而超出了经典\Term{algebraic-geometry}{代数几何}的范围。

\textbf{5.} 我们认为，没有必要在每讲书面讲义的开头注明相应口头报告的日期。这里只需说明，油印讲义（第 I 至 XXVI 讲）的次序确实与口头报告的次序一致。此外，定稿的撰写时间有时显著晚于口头报告，而且内容往往有相当大的不同；书面文本一般更详尽、更完整（例如第 IV 和 VII\textsubscript{B} 讲），甚至明显更一般（例如第 XII 或 VII\textsubscript{B} 讲）。另有一些书面讲义并不对应任何口头报告（第 VI\textsubscript{B}、VII\textsubscript{A}、XV、XVI、XVII 讲，以及第 XXVI 讲的大部分），将它们撰写并插入油印本，有的是为了给其他各讲提供方便的参考（第 VI\textsubscript{B} 讲尤其如此），有的是因为它们是已发展的概念和技术的自然延伸。因此，学习本讨论班的其余部分，尤其学习其中关于\Term{reductive-group-scheme}{约化群概形}的部分，并不需要阅读第 VII\textsubscript{A}、VII\textsubscript{B}、XV、XVI、XVII 讲。

\textbf{6.} 在\Term{scheme}{概形}论中，我们主要使用 EGA I 所述的\Term{scheme}{概形}一般语言，SGA 1，第 I 至 V 讲所述的\Term{flat-morphism}{平坦态射}、\Term{etale}{平展态射}、\Term{smooth-morphism}{光滑态射}的概念，以及 SGA 1，第 VIII 讲的\Term{faithfully-flat-descent}{忠实平坦下降}理论。我们尽可能避免提出无用的 Noether 性假设；这却迫使我们用“\Term{finite-presentation}{有限表示}”假设代替惯用的“\Term{finite-type}{有限型}”假设。关于\Term{finite-presentation}{有限表示}\Term{morphism}{态射}的概念，读者可查阅 EGA IV，1.4 和 1.6。SGA 1，第 I 至 IV 讲中的结果大多是在 Noether 条件下陈述的；EGA IV 将在一般情形中发展这些结果\footnote{\textsuperscript{(****)} 自本引言写成以来，EGA IV 的四个部分（第 1 至 21 节）已全部出版。}，并将详细阐述把某些类型的陈述（涉及\Term{finite-presentation}{有限表示}假设）归结到 Noether 情形的标准方法（EGA IV，第 8、9、11 节）。不愿承认 EGA IV 这些结果的读者，可以通过假设底\Term{prescheme}{预概形}\Term{locally-noetherian}{局部 Noether} 来简化某些陈述或其证明。不过，当证明是从 $\widehat A$ \Term{descent}{下降}到 $A$，其中 $\widehat A$ 为 \Term{noetherian-local-ring}{Noether 局部环} $A$ 的\Term{completion}{完备化}时，
\SourcePage{6}
这种做法会遇到困难，因为这种方法会引入一般并非 Noether 的环 $\widehat A\otimes_A\widehat A$。

\textbf{7.} 引用采用通常的十进制编号方式：5.7.11 指同一讲中该编号的命题（或\Term{lemma}{引理}、定义等）；XVII 7.8 中的罗马数字表示讲次。标准文献引用使用下列缩写：
\begin{description}
\item[Bible] 1956/58 年 Chevalley 讨论班“Groupes de Lie algébriques”。
\item[EGA X, x.y.z] J. Dieudonné 与 A. Grothendieck，\emph{Éléments de Géométrie Algébrique}，第 X 章，陈述 x.y.z（或小节 x.y 等）。
\item[SGA n, X y.z] Bois-Marie \Term{algebraic-geometry}{代数几何}讨论班，第 $n$ 年，第 X 讲的陈述 y.z。
\item[TDTE] A. Grothendieck，\emph{Techniques de descente et Théorèmes d'existence en Géométrie Algébrique}，1959 至 1962 年间在 Bourbaki 讨论班所作的报告。
\end{description}

\textbf{编者增补注 (1)。} 自 A. Grothendieck 写下脚注 (**) 以来，\Term{reductive-group}{约化群}的商问题已有很大发展。对于任意域 $k$ 上的\Term{affine-algebraic-group}{仿射代数群}，由 Mumford 引入的“正确”概念是“\Term{geometrically-reduced}{几何约化}群”。（另一方面，若 $G$ 的每个\Term{rational-representation}{有理表示}都完全可约，就称 $G$“\Term{linearly-reductive}{线性约化}”；但如脚注 (**) 所述，当 $\operatorname{char}(k)>0$ 时，这个条件\Term{restriction}{限制}过强。）由 M. Nagata 和 W. J. Haboush 的结果（[Na64]、[NM64]、[Ha75]），\Term{geometrically-reduced}{几何约化}的 $k$-群正好就是本讨论班意义下的\Term{reduced}{约化} $k$-群。这方面可参见 Mumford 著作的后续版本（[MF82]）。此外，C. S. Seshadri（[Se77]）将“\Term{geometrically-reduced}{几何约化}性”的概念以及它对于取商的推论推广到了任意底的情形；J.-L. Colliot-Thélène 与 J.-J. Sansuc 的文章 [CTS79] 中还载有 M. Raynaud 的补充。最后，关于对\Term{algebraic-group}{代数群}，或更一般地对\Term{groupoid}{群胚}（参见本讨论班第 V 讲）取商的较新发展，我们推荐 J. Kollár 的文章 [Ko97]，以及 S. Keel 与 S. Mori 的文章 [KM97]。

\textbf{参考文献}
\begin{description}
\item[{[CTS79]}] J.-L. Colliot-Thélène \& J.-J. Sansuc，\emph{Fibrés quadratiques et composantes connexes réelles}，Math. Ann. \textbf{244}（1979），105--134。
\item[{[Ha75]}] W. J. Haboush，\emph{Reductive groups are geometrically reductive}，Ann. of Math. \textbf{102}（1975），第 1 期，67--83。
\item[{[KM97]}] S. Keel \& S. Mori，\emph{Quotient by groupoids}，Ann. of Math. \textbf{145}（1997），第 1 期，193--213。
\item[{[Ko97]}] J. Kollár，\emph{Quotient spaces modulo algebraic groups}，Ann. of Math. \textbf{145}（1997），第 1 期，33--79。
\item[{[MF82]}] D. Mumford \& J. Fogarty，\emph{Geometric invariant theory}，第 2 版，Springer-Verlag，1982；（第 3 版与 F. Kirwan 合著，1994）。
\item[{[Na64]}] M. Nagata，\emph{Invariants of a group in an affine ring}，J. Math. Kyoto Univ. \textbf{3}（1964），第 3 期，369--377。
\item[{[NM64]}] M. Nagata \& T. Miyata，\emph{Note on semi-reductive groups}，J. Math. Kyoto Univ. \textbf{3}（1964），第 3 期，379--382。
\item[{[Se77]}] C. S. Seshadri，\emph{Geometric reductivity over an arbitrary base}，Adv. Math. \textbf{26}（1977），第 3 期，225--274。
\end{description}

\SourcePage{7}
\SGAchapter{I}{代数结构；群上同调}{Structures algébriques. Cohomologie des groupes}
\begin{center}M. Demazure\end{center}
本讲分为两部分：第一部分汇集一些一般定义，并规定后文经常使用的记号；第二部分讨论\Term{group-cohomology}{群上同调}，最终得到\Term{theorem}{定理} 5.3.3（\Term{diagonalizable-group}{可对角化群}的\Term{cohomology}{上同调}消失）。

我们一次选定一个\Term{universe}{宇宙}\footnote{\textsuperscript{(1)} 编者注：参见 SGA 4，第 I 讲，第 0 节及附录；也可参见 [DG70]，第 xxvi 页的讨论。}。所有定义和构造都相对于这个\Term{universe}{宇宙}。我们将系统地容许下列语言上的滥用：为定义\Term{functor}{函子} $f:\mathcal C\to\mathcal C'$，只要对于 $\mathcal C$ 中\Term{arrow}{箭头} $h$ 应如何定义 $f(h)$ 没有歧义，就只定义每个\Term{object}{对象} $S$ 对应的 $\mathcal C'$ 中的\Term{object}{对象} $f(S)$。实际叙述中，我们将说：设 $f:\mathcal C\to\mathcal C'$ 为由 $f(S)=\cdots$ 定义的\Term{functor}{函子}。

\SGAsection{1}{一般事项}
\textbf{1.1.} 设 $\mathcal C$ 为一个\Term{category}{范畴}。用 $\widehat{\mathcal C}$ 表示从 $\mathcal C$ 到\Term{set}{集合}\Term{category}{范畴} $(\mathrm{Ens})$ 的\Term{contravariant-functor}{反变函子}\Term{category}{范畴} $\operatorname{Hom}(\mathcal C^\circ,(\mathrm{Ens}))$\footnote{\textsuperscript{(2)} 编者注：它称为 $\mathcal C$ 上的\Term{presheaf}{预层}\Term{category}{范畴}，参见 IV 4.3.1。}。有一个\Term{canonical-functor}{典范函子} $h:\mathcal C\to\widehat{\mathcal C}$，它将每个 $X\in\operatorname{Ob}(\mathcal C)$ 映到\Term{functor}{函子} $h_X$，其中
\[
h_X(S)=\operatorname{Hom}(S,X).
\]
对于每个 $F\in\operatorname{Ob}(\widehat{\mathcal C})$，可定义（例如参见 EGA $0_{\mathrm{III}}$，8.1.4）一个\Term{bijection}{双射}
\[
\operatorname{Hom}(h_X,F)\xrightarrow{\sim}F(X).
\]
\footnote{\textsuperscript{(3)} 编者注：这个结果通常称为“Yoneda \Term{lemma}{引理}”；我们在其他编者注中将使用这个名称。}
特别地，对于 $\mathcal C$ 中任意两个\Term{object}{对象} $X,Y$，下列\Term{canonical}{典范}映射是\Term{bijection}{双射}：
\[
\operatorname{Hom}_{\mathcal C}(X,Y)\xrightarrow{\sim}
\operatorname{Hom}_{\widehat{\mathcal C}}(h_X,h_Y);
\]
\SourcePage{8}
也就是说，\Term{functor}{函子} $h$ 是\Term{fully-faithful}{全忠实}的。因此，它给出 $\mathcal C$ 到 $\widehat{\mathcal C}$ 的某个\Term{full-subcategory}{满子范畴}的\Term{isomorphism}{同构}，并给出 $\mathcal C$ 与 $\widehat{\mathcal C}$ 中由\Term{representable-functor}{可表函子}（即\Term{isomorphism}{同构}于形如 $h_X$ 的\Term{functor}{函子}的那些\Term{functor}{函子}）组成的\Term{full-subcategory}{满子范畴}之间的\Term{equivalence}{等价}。以后我们常将 $X$ 与 $h_X$ 等同。以下各款的目的，就是说明这种等同不会造成问题。

\begin{SGAresult}{注记}{1.1.1}
\footnote{\textsuperscript{(4)} 编者注：增补了本注记。}
有时需要下面的变式。设 $\mathcal D$ 是 $\mathcal C$ 的\Term{full-subcategory}{满子范畴}，且 $X,Y\in\operatorname{Ob}(\mathcal D)$；用 $h'_X,h'_Y$ 表示 $h_X,h_Y$ 在 $\mathcal D$ 上的\Term{restriction}{限制}。那么
\[
\operatorname{Hom}_{\mathcal C}(X,Y)=\operatorname{Hom}_{\mathcal D}(X,Y)
=\operatorname{Hom}_{\widehat{\mathcal D}}(h'_X,h'_Y),
\]
因而：给定一个\Term{morphism}{态射} $X\to Y$，“等同于”对每个 $T\in\operatorname{Ob}(\mathcal D)$ 给定一个关于 $T$ \Term{functorial}{函子性}的映射 $\phi(T):\operatorname{Hom}(T,X)\to\operatorname{Hom}(T,Y)$。
\end{SGAresult}

\textbf{1.2.}\footnote{\textsuperscript{(5)} 编者注：调整了叙述次序，以便先引入\Term{fiber-product}{纤维积}，再引入\Term{monomorphism}{单态射}；参见编者注 (9)。}
若对于每个 $S$，$F(S)$ 是 $G(S)$ 的\Term{subset}{子集}，就称 $F$ 为 $G$ 的\Term{subobject}{子对象}（或\Term{subfunctor}{子函子}）。

$\widehat{\mathcal C}$ 中存在“任意”的\Term{projective-limit}{投射极限}，其计算公式为
\[
(\varprojlim_i F_i)(S)=\varprojlim_i F_i(S).
\]
\footnote{\textsuperscript{(6)} 编者注：同样，“任意”的\Term{inductive-limit}{归纳极限}存在，并“逐自变量”计算，即 $(\varinjlim_iF_i)(S)=\varinjlim_iF_i(S)$；但一般而言，\Term{functor}{函子} $h$ 不与\Term{inductive-limit}{归纳极限}交换。}
特别地，\Term{fiber-product}{纤维积}由下式定义：
\[
(F\times_G F')(S)=F(S)\times_{G(S)}F'(S).
\]
\footnote{\textsuperscript{(7)} 编者注：特别地，一对\Term{morphism}{态射} $u,v:F\rightrightarrows G$ 的核是 $F$ 的\Term{subfunctor}{子函子} $\operatorname{Ker}(u,v)$，由 $\operatorname{Ker}(u,v)(S)=\{x\in F(S)\mid u(x)=v(x)\}$ 定义。}
我们选取 $e(S)=\{\varnothing\}$ 的\Term{functor}{函子} $e$ 作为 $\widehat{\mathcal C}$ 的\Term{final-object}{终对象}\footnote{\textsuperscript{(8)} 编者注：$\{\varnothing\}$（空集的幂集）表示单元素\Term{set}{集合}。}。每个 $F\in\operatorname{Ob}(\widehat{\mathcal C})$ 都有唯一\Term{morphism}{态射}到 $e$，并置
\[
F\times F'=F\times_eF'.
\]
\Term{functor}{函子} $h$ 与\Term{inverse-limit}{投射极限}交换；特别地，对于 $X,X'\in\operatorname{Ob}(\mathcal C)$，$X\times X'$ 存在（分别地，$\mathcal C$ 有\Term{final-object}{终对象} $e$）的充分必要条件是 $h_X\times h_{X'}$ \Term{representable}{可表}（分别地，$e$ \Term{representable}{可表}），并且
\[
h_X\times h_{X'}\simeq h_{X\times X'},\qquad h_e\simeq e.
\]
\SourcePage{9}
$\widehat{\mathcal C}$ 中的\Term{monomorphism}{单态射}，恰是那些对于每个 $S\in\operatorname{Ob}(\mathcal C)$，相应的\Term{set}{集合}映射 $F(S)\to G(S)$ 均为\Term{injectivity}{单射}的\Term{morphism}{态射} $F\to G$\footnote{\textsuperscript{(9)} 编者注：若每个 $F(S)\to G(S)$ 都是\Term{injectivity}{单射}，显然 $F\to G$ 是\Term{monomorphism}{单态射}；反命题可通过考虑图 $F\times_GF\rightrightarrows F\to G$ 看出。由此得到：“$F\to G$ 是\Term{monomorphism}{单态射}，当且仅当\Term{diagonal-morphism}{对角态射} $F\to F\times_GF$ 为\Term{isomorphism}{同构}”（参见 EGA I，5.3.8）。同样，若每个 $F(S)\to G(S)$ 都是\Term{surjectivity}{满射}，则显然 $F\to G$ 是\Term{epimorphism}{满态射}；反命题可通过考虑\Term{pushout}{并合和} $G\amalg_FG$ 看出，参见第 IV 讲\Term{lemma}{引理} 4.4.4 的证明。}。

\textbf{\Term{functor}{函子} $\Gamma$。} 对于每个 $F\in\operatorname{Ob}(\widehat{\mathcal C})$，置
\[
\Gamma(F)=\operatorname{Hom}(e,F);
\]
因此，$\Gamma(F)$ 的一个元素就是一个族 $(\gamma_S)_{S\in\operatorname{Ob}(\mathcal C)}$，其中 $\gamma_S\in F(S)$，并且对于 $\mathcal C$ 中的每个\Term{arrow}{箭头} $f:S'\to S''$ 都有 $F(f)(\gamma_{S''})=\gamma_{S'}$。

对于 $X\in\operatorname{Ob}(\mathcal C)$，置 $\Gamma(X)=\Gamma(h_X)$。若 $\mathcal C$ 有\Term{final-object}{终对象} $e$，便有\Term{isomorphism}{同构} $\Gamma(X)\simeq\operatorname{Hom}(e,X)$。

\textbf{1.3.} 设 $S\in\operatorname{Ob}(\mathcal C)$。用 $\mathcal C_{/S}$ 表示 $\mathcal C$ 中位于 $S$ 之上的\Term{object}{对象}的\Term{category}{范畴}，即其\Term{object}{对象}为 $\mathcal C$ 中的\Term{arrow}{箭头} $f:T\to S$，而 $\operatorname{Hom}(f,f')$ 是 $\operatorname{Hom}(T,T')$ 中满足 $f=f'\circ u$ 的 $u$ 所组成的\Term{subset}{子集}。若 $\mathcal C$ 有\Term{final-object}{终对象} $e$，则 $\mathcal C_{/e}$ \Term{isomorphism}{同构}于 $\mathcal C$。\Term{category}{范畴} $\mathcal C_{/S}$ 有\Term{final-object}{终对象}，即恒等\Term{arrow}{箭头} $S\to S$。

若 $f:T\to S$ 为 $\mathcal C_{/S}$ 的\Term{object}{对象}，则\Term{constructible}{可构造}\Term{category}{范畴} $(\mathcal C_{/S})_{/f}$；我们滥用记号，将它记为 $(\mathcal C_{/S})_{/T}$，并有\Term{canonical-isomorphism}{典范同构}
\[
\mathcal C_{/T}\simeq(\mathcal C_{/S})_{/T}.
\]
这一构造也适用于\Term{category}{范畴} $\widehat{\mathcal C}$；特别地，定义了\Term{category}{范畴} $\widehat{\mathcal C}_{/h_S}$。另一方面，也\Term{constructible}{可构造}\Term{category}{范畴} $\widehat{\mathcal C_{/S}}$。

若 $f:T\to S$ 为 $\mathcal C_{/S}$ 的\Term{object}{对象}，则 $\Gamma(f)$ 可等同于 $T$ 在 $S$ 上的\Term{section}{截面}的\Term{set}{集合} $\Gamma(T/S)$，即 $f$ 的右逆\Term{arrow}{箭头} $S\to T$ 所成的\Term{set}{集合}。注意，$h_f:h_T\to h_S$ 是 $\widehat{\mathcal C}_{/h_S}$ 的\Term{object}{对象}，并且
\[
\Gamma(h_f)\simeq\Gamma(h_T/h_S)\simeq\Gamma(T/S)\simeq\Gamma(f).
\]

\textbf{1.4.} 现在要定义\Term{category}{范畴} $\widehat{\mathcal C_{/S}}$ 与 $\widehat{\mathcal C}_{/h_S}$ 的一个\Term{equivalence-of-categories}{等价}，亦即证明：“给定 $\mathcal C$ 中位于 $S$ 之上的\Term{object}{对象}的\Term{category}{范畴}上的一个\Term{functor}{函子}，等同于给定 $\mathcal C$ 上的一个\Term{functor}{函子}，并附以一个到 $h_S$ 的\Term{morphism}{态射}。”

\textbf{(i) 构造 $\alpha_S:\widehat{\mathcal C}_{/h_S}\to\widehat{\mathcal C_{/S}}$。}
首先，设 $H:F\to h_S$ 是 $\widehat{\mathcal C}_{/h_S}$ 的\Term{object}{对象}。我们要定义 $\mathcal C_{/S}$ 上的\Term{functor}{函子} $\alpha_S(H)$。先取 $\mathcal C_{/S}$ 的\Term{object}{对象} $f:T\to S$；将 $\alpha_S(H)(f)$ 定义为元素 $f\in h_S(T)$ 在映射 $H(T):F(T)\to h_S(T)$ 下的原像\footnote{\textsuperscript{(10)} 编者注：例如，若 $F=h_X$，则 $H$ 对应于一个\Term{morphism}{态射} $h:X\to S$，而 $H(T):\operatorname{Hom}_{\mathcal C}(T,X)\to\operatorname{Hom}_{\mathcal C}(T,S)$ 为映射 $g\mapsto h\circ g$，故 $\alpha_S(h_X)=\operatorname{Hom}_{\mathcal C_{/S}}(-,X)$。}。
\SourcePage{10}
再设 $u:f\to f'$ 是 $\mathcal C_{/S}$ 中的\Term{arrow}{箭头}，则 $F(u):F(T')\to F(T)$ 诱导一个从 $\alpha_S(H)(f')$ 到 $\alpha_S(H)(f)$ 的映射，记为 $\alpha_S(H)(u)$。立刻可核知，映射
\[
f\longmapsto\alpha_S(H)(f),\qquad u\longmapsto\alpha_S(H)(u)
\]
确实定义了 $\mathcal C_{/S}$ 上的\Term{functor}{函子}，因而定义了 $\widehat{\mathcal C_{/S}}$ 的\Term{object}{对象} $\alpha_S(H)$。

最后，设 $H:F\to h_S$ 与 $H':F'\to h_S$ 是 $\widehat{\mathcal C}_{/h_S}$ 的两个\Term{object}{对象}，$U:H\to H'$ 是其中的一个\Term{morphism}{态射}：
\[
\begin{tikzcd}
F\arrow[rr,"U"]\arrow[dr,"H"']&&F'\arrow[dl,"H'"]\\
&h_S&
\end{tikzcd}
\]
那么，对于每个 $f:T\to S$，映射 $U(T):F(T)\to F'(T)$ 诱导映射
\[
\alpha_S(U)(f):\alpha_S(H)(f)\longrightarrow\alpha_S(H')(f),
\]
从而定义了\Term{morphism-of-functors}{函子态射}
\[
\alpha_S(U):\alpha_S(H)\longrightarrow\alpha_S(H').
\]
容易核知，映射
\[
H\longmapsto\alpha_S(H),\qquad U\longmapsto\alpha_S(U)
\]
确实定义了\Term{functor}{函子} $\alpha_S:\widehat{\mathcal C}_{/h_S}\to\widehat{\mathcal C_{/S}}$。

\textbf{(ii)}
\begin{SGAresult}{命题}{1.4.1}
\Term{functor}{函子} $\alpha_S:\widehat{\mathcal C}_{/h_S}\to\widehat{\mathcal C_{/S}}$ 是\Term{category-equivalence}{范畴等价}。
\end{SGAresult}
这里只说明构造一个\Term{quasi-inverse}{拟逆}\Term{functor}{函子} $\beta_S:\widehat{\mathcal C_{/S}}\to\widehat{\mathcal C}_{/h_S}$ 的原则。设 $G$ 是 $\mathcal C_{/S}$ 上的\Term{functor}{函子}；对于 $\mathcal C$ 中每个\Term{object}{对象} $T$，置
\[
\beta_S(G)(T)=\coprod_{f\in\operatorname{Hom}(T,S)=h_S(T)}G(f),
\]
即对所有 $f\in\operatorname{Hom}(T,S)=h_S(T)$ 取\Term{set}{集合} $G(f)$ 的和。这定义了 $\mathcal C$ 上的\Term{functor}{函子} $\beta_S(G)$，并且它带有显然的到 $h_S$ 的\Term{projection}{投影}。

\textbf{1.5.} \Term{equivalence}{等价} $\alpha_S$ 与\Term{functor}{函子} $\Gamma$ 交换。换言之，若 $H:F\to h_S$ 是 $\widehat{\mathcal C}_{/h_S}$ 的\Term{object}{对象}，$\alpha_S(H)$ 是相应的 $\widehat{\mathcal C_{/S}}$ 的\Term{object}{对象}，则
\[
\Gamma(\alpha_S(H))\simeq\Gamma(H)\simeq\Gamma(F/h_S).
\]
\Term{equivalence}{等价} $\alpha_S$ 与\Term{functor}{函子} $h$ 交换，即若 $f:T\to S$ 为 $\mathcal C_{/S}$ 的\Term{object}{对象}，则 $h_f:h_T\to h_S$ 为 $\widehat{\mathcal C}_{/h_S}$ 的\Term{object}{对象}，它在 $\alpha_S$ 下的像正是 $h_{\mathcal C_{/S}}(f)$，其中
\[
h_{\mathcal C_{/S}}:\mathcal C_{/S}\longrightarrow\widehat{\mathcal C_{/S}}
\]
是\Term{canonical-functor}{典范函子}\footnote{\textsuperscript{(11)} 编者注：参见编者注 (10)。}。因此：
\SourcePage{11}
\begin{SGAresult}{命题}{1.5.1}
设 $H:F\to h_S$ 为 $\widehat{\mathcal C}_{/h_S}$ 的\Term{object}{对象}。\Term{functor}{函子} $\alpha_S(H):(\mathcal C_{/S})^\circ\to(\mathrm{Ens})$ \Term{representable}{可表}，当且仅当 $F:\mathcal C^\circ\to(\mathrm{Ens})$ \Term{representable}{可表}；若 $F\simeq h_T$，则 $\alpha_S(H)$ 由 $\mathcal C_{/S}$ 的\Term{object}{对象} $T\to S$ 代表。
\end{SGAresult}
\Term{equivalence}{等价} $\alpha_S$ 对 $S$ 是传递的：若 $f:T\to S$ 为 $\mathcal C_{/S}$ 的\Term{object}{对象}，则有由\Term{equivalence}{等价}组成的\Term{commutative-diagram}{交换图}
\[
\begin{tikzcd}[column sep=large]
(\widehat{\mathcal C}_{/h_S})_{/h_T}
 \arrow[r,"\alpha_S/h_T"]\arrow[d,"\sim"']&
(\widehat{\mathcal C_{/S}})_{/h_T}\arrow[r,"\alpha_f"]&
\widehat{(\mathcal C_{/S})_{/T}}\arrow[d,"\sim"]\\
\widehat{\mathcal C}_{/h_T}\arrow[rr,"\alpha_T"]&&\widehat{\mathcal C_{/T}}
\end{tikzcd}
\]
其中 $\alpha_S/h_T$（暂时）表示将\Term{functor}{函子} $\alpha_S$ \Term{restriction}{限制}在位于 $h_T$ 之上的\Term{object}{对象}上（参见 1.6）。

\textbf{1.6. \Term{functor}{函子}的\Term{base-change}{基变换}。} 对每个 $S\in\operatorname{Ob}(\mathcal C)$，有\Term{canonical-functor}{典范函子}
\[
i_S:\mathcal C_{/S}\longrightarrow\mathcal C,
\]
若 $f$ 为\Term{arrow}{箭头} $T\to S$，则定义 $i_S(f)=T$。若 $f:T\to S$ 是 $\mathcal C_{/S}$ 的\Term{object}{对象}，则用 $i_{T/S}=i_f$ 表示\Term{functor}{函子}
\[
i_{T/S}:(\mathcal C_{/S})_{/T}\longrightarrow\mathcal C_{/S},
\]
并有\Term{commutative-diagram}{交换图}
\[
\begin{tikzcd}
(\mathcal C_{/S})_{/T}\arrow[r,"i_{T/S}"]\arrow[dr,"\sim"']&\mathcal C_{/S}\arrow[r,"i_S"]&\mathcal C\\
&\mathcal C_{/T}\arrow[ur,"i_T"']&
\end{tikzcd}
\]
也就是说，像今后那样将 $(\mathcal C_{/S})_{/T}$ 与 $\mathcal C_{/T}$ 等同，就有
\[
i_S\circ i_{T/S}=i_T.
\]
同样，当 $\mathcal C$ 有\Term{final-object}{终对象} $e$，并将 $\mathcal C$ 与 $\mathcal C_{/e}$ 等同时，$i_{S/e}:\mathcal C_{/S}\to\mathcal C_{/e}$ 就等同于 $i_S$。

对于 $X\in\operatorname{Ob}(\mathcal C)$（分别地，$Y\in\operatorname{Ob}(\mathcal C_{/S})$），用 $p_S(X)$（分别地，$p_{T/S}(Y)$）表示由 $X\times S$（分别地，$Y\times_ST$）及其第二\Term{projection}{投影}所定义的 $\mathcal C_{/S}$（分别地，$\mathcal C_{/T}$）的\Term{object}{对象}，只要它存在：
\[
\begin{tikzcd}
X\times S\arrow[d,"p_S(X)"']&&Y\times_ST\arrow[d,"p_{T/S}(Y)"]\\
S&&T.
\end{tikzcd}
\]
\SourcePage{12}
这个（部分定义的）\Term{functor}{函子} $p_S$（分别地，$p_{T/S}$）称为\Term{base-change}{基变换函子}。按照积（分别地，\Term{fiber-product}{纤维积}）的定义，它是\Term{functor}{函子} $i_S$（分别地，$i_{T/S}$）的\Term{right-adjoint}{右伴随}\Term{functor}{函子}\footnote{\textsuperscript{(12)} 编者注：即若 $g:U\to S$（分别地，$h:V\to T$）为 $\mathcal C_{/S}$（分别地，$\mathcal C_{/T}$）的\Term{object}{对象}，则 $i_S(g)=U$，而 $i_{T/S}(h)$ 为 $\mathcal C_{/S}$ 的\Term{object}{对象} $f\circ h:V\to S$，并且
\[
\operatorname{Hom}_{\mathcal C_{/S}}(U,X\times S)\simeq\operatorname{Hom}_{\mathcal C}(U,X),
\]
分别地，
\[
\operatorname{Hom}_{\mathcal C_{/T}}(V,X\times_ST)\simeq\operatorname{Hom}_{\mathcal C_{/S}}(V,X).
\]
}。也记作
\[
p_S(X)=X_S,\qquad p_{T/S}(Y)=Y_T.
\]
\Term{functor}{函子} $i_S$ 定义一个（\Term{restriction}{限制}）\Term{functor}{函子}
\[
i_S^*:\widehat{\mathcal C}\longrightarrow\widehat{\mathcal C_{/S}};
\]
记 $F_S=i_S^*(F)=F\circ i_S$。显然有
\[
i_{T/S}^*\circ i_S^*=i_T^*,
\]
即对于每个\Term{functor}{函子} $F\in\operatorname{Ob}(\widehat{\mathcal C})$，
\[
(F_S)_T=F_T.
\]
这个记号需要以下说明。
\begin{SGAresult}{命题}{1.6.1}
\Term{functor}{函子} $(h_X)_S:(\mathcal C_{/S})^\circ\to(\mathrm{Ens})$ \Term{representable}{可表}的充分必要条件是积 $X\times S$ 存在。这时
\[
(h_X)_S\simeq h_{X_S}.
\]
\end{SGAresult}
这表明，$F_S$ 有两种解释：\Term{functor}{函子} $F$ 在 $\mathcal C_{/S}$ 上的\Term{restriction}{限制}，以及由\Term{base-change}{基变换} $e\leftarrow S$ 得到的\Term{functor}{函子}。于是采用如下记法：
\[
\begin{tikzcd}
F&F_S&F_T\\
e&S\arrow[l]&T\arrow[l]
\end{tikzcd}
\]
它恰好反映了上述两种解释。

注意，我们有
\[
\Gamma(F_S)\simeq\operatorname{Hom}(h_S,F)\simeq F(S),
\]
特别地，
\[
\Gamma(X_S)\simeq\operatorname{Hom}(S,X).
\]

\textbf{1.7.0.}\footnote{\textsuperscript{(13)} 编者注：增补了下面的\Term{lemma}{引理}（参见 SGA 4，I.3.4），它用于 1.7.1 的证明，并将在后文多次用到。}
设 $E$ 是 $\widehat{\mathcal C}$ 的\Term{object}{对象}。考虑 $\mathcal C$ 中位于 $E$ 之上的\Term{object}{对象}的\Term{category}{范畴} $\mathcal C_{/E}$：其\Term{object}{对象}为二元组 $(V,\rho)$，其中 $V$ 是 $\mathcal C$ 的\Term{object}{对象}，$\rho:h_V\to E$ 是 $\widehat{\mathcal C}$ 中的\Term{morphism}{态射}，即 $\rho\in E(V)$；从 $(V,\rho)$ 到 $(V',\rho')$ 的\Term{morphism}{态射}，
\SourcePage{13}
是满足 $\rho'\circ f=\rho$（即 $E(f)(\rho')=\rho$）的一个 $\mathcal C$-\Term{morphism}{态射} $f:V\to V'$。用 $L$ 表示\Term{functor}{函子}
\[
\varinjlim_{(V,\rho)\in\mathcal C_{/E}}h_V,
\]
即对于每个 $S\in\operatorname{Ob}(\mathcal C)$，$L(S)=\varinjlim_{(V,\rho)}h_V(S)$ 是三元组 $(V,\rho,v)$ 的\Term{equivalence-class}{等价类}所成的\Term{set}{集合}，其中 $v:S\to V$ 是 $\mathcal C$-\Term{morphism}{态射}，并且对每个满足 $\rho'\circ f=\rho$ 的 $\mathcal C$-\Term{morphism}{态射} $f:V\to V'$，将 $(V,\rho,v)$ 与 $(V',\rho',f\circ v)$ 等同。

将 $(V,\rho,v)$ 的类映为 $E(S)$ 中的元素 $\rho\circ v$，就得到良定义的映射 $\phi_E(S)$，它定义了\Term{morphism-of-functors}{函子态射}
\[
\phi_E:\varinjlim_{(V,\rho)\in\mathcal C_{/E}}h_V\longrightarrow E.
\]
\begin{SGAresult}{\Term{lemma}{引理}}{}
$\phi_E$ 是\Term{isomorphism}{同构}。
\end{SGAresult}
\begin{proof}
设 $S\in\operatorname{Ob}(\mathcal C)$。每个 $x\in E(S)$ 都是三元组 $(S,x,\operatorname{id}_S)$ 在 $\phi_E(S)$ 下的像；所以 $\phi_E(S)$ 是\Term{surjectivity}{满射}。另一方面，设 $\ell_1=(V_1,\rho_1,v_1)$ 和 $\ell_2=(V_2,\rho_2,v_2)$ 是 $L(S)$ 中在 $E(S)$ 内有相同像的两个元素；置 $\gamma=\rho_1\circ v_1=\rho_2\circ v_2$。那么，在 $L(S)$ 中，$\ell_1$ 和 $\ell_2$ 都等于三元组 $(S,\gamma,\operatorname{id}_S)$ 的类。这表明 $\phi_E(S)$ 是\Term{injectivity}{单射}。
\end{proof}
\begin{SGAresult}{推论}{}
对于 $\widehat{\mathcal C}$ 的每个\Term{object}{对象} $F$，有
\[
\operatorname{Hom}(E,F)=\varprojlim_{(V,\rho)\in\mathcal C_{/E}}F(V).
\]
\end{SGAresult}

\textbf{1.7. $\underline{\operatorname{Hom}}$、$\underline{\operatorname{Isom}}$ 等\Term{object}{对象}。}
设 $F,G$ 是 $\widehat{\mathcal C}$ 的两个\Term{object}{对象}。用下列方式定义 $\widehat{\mathcal C}$ 的另一个\Term{object}{对象}：
\[
\begin{aligned}
\underline{\operatorname{Hom}}(F,G)(S)
&=\operatorname{Hom}_{\widehat{\mathcal C_{/S}}}(F_S,G_S)\\
&\simeq\operatorname{Hom}_{\widehat{\mathcal C}_{/h_S}}(F\times h_S,G\times h_S)
\simeq\operatorname{Hom}_{\widehat{\mathcal C}}(F\times h_S,G).
\end{aligned}
\]
这样定义的\Term{object}{对象} $\underline{\operatorname{Hom}}(F,G)$ 具有下列性质：

(i) $\underline{\operatorname{Hom}}(e,G)\simeq G$\footnote{\textsuperscript{(14)} 编者注：而且，若 $E$ 是 $\widehat{\mathcal C}$ 的第三个\Term{object}{对象}，则有 $\underline{\operatorname{Hom}}(E,F\times G)\simeq\underline{\operatorname{Hom}}(E,F)\times\underline{\operatorname{Hom}}(E,G)$。}。

(ii) $\underline{\operatorname{Hom}}$ 的形成与\Term{base-change}{基扩张}交换：
\[
\underline{\operatorname{Hom}}(F_S,G_S)\simeq\underline{\operatorname{Hom}}(F,G)_S.
\]
(iii) $(F,G)\mapsto\underline{\operatorname{Hom}}(F,G)$ 是一个\Term{bifunctor}{双函子}，对 $F$ \Term{contravariant}{反变}，对 $G$ \Term{covariant}{协变}。

这三条性质由定义即明。

我们将证明，对于 $\widehat{\mathcal C}$ 的每个\Term{object}{对象} $E$，有
\[
\operatorname{Hom}(E,\underline{\operatorname{Hom}}(F,G))\simeq\operatorname{Hom}(E\times F,G).
\]
\SourcePage{14}
设 $\phi:E\times F\to G$；我们要给它对应一个从 $E$ 到 $\underline{\operatorname{Hom}}(F,G)$ 的\Term{morphism}{态射}。设 $S'\to S$ 是 $\mathcal C$ 中的\Term{arrow}{箭头}。有映射
\[
E(S)\times F(S')\longrightarrow E(S')\times F(S')\xrightarrow{\phi(S')}G(S').
\]
因此，$E(S)$ 的每个元素 $e$ 对每个 $S'\to S$ 都定义一个关于 $S'$ \Term{functorial}{函子性}的映射 $F(S')\to G(S')$，即 $\underline{\operatorname{Hom}}(F,G)(S)$ 的一个元素 $\theta_\phi(e)$。于是得到一个“关于 $E$ \Term{functorial}{函子性}”的映射
\[
\phi\longmapsto\theta_\phi,\qquad
\operatorname{Hom}(E\times F,G)\longrightarrow\operatorname{Hom}(E,\underline{\operatorname{Hom}}(F,G)).
\]
\begin{SGAresult}{命题}{1.7.1}
\footnote{\textsuperscript{(15)} 编者注：增补了 (b)，它将在 II.1 和 II.3.11 中用到。此外，还勾勒了 (a) 的第二个更直接的证明。}
设 $E,F,G\in\operatorname{Ob}(\widehat{\mathcal C})$。

(a) 映射 $\phi\mapsto\theta_\phi$ 是\Term{bijection}{双射}：
\[
\operatorname{Hom}(E\times F,G)\xrightarrow{\sim}\operatorname{Hom}(E,\underline{\operatorname{Hom}}(F,G)).
\]
(b) 此外，有\Term{functor}{函子}\Term{isomorphism}{同构}：
\[
\underline{\operatorname{Hom}}(E,\underline{\operatorname{Hom}}(F,G))\simeq\underline{\operatorname{Hom}}(E\times F,G).
\]
\end{SGAresult}
\begin{proof}
(a) 将两端看成 $E$ 的\Term{functor}{函子}。当 $E=h_X$ 时，所述结果成立；事实上，此时它正是\Term{functor}{函子} $\underline{\operatorname{Hom}}(F,G)$ 的定义。另一方面，两端作为 $E$ 的\Term{functor}{函子}，都把\Term{inductive-limit}{归纳极限}变为\Term{inverse-limit}{投射极限}。最后，由\Term{lemma}{引理} 1.7.0，$\widehat{\mathcal C}$ 的每个\Term{object}{对象} $E$ 都\Term{isomorphism}{同构}于 $h_X$ 的\Term{inductive-limit}{归纳极限}，其中 $X$ 遍历\Term{category}{范畴} $\mathcal C_{/E}$。这就证明了 (a)。

\textsuperscript{(15)} 再勾勒 (a) 的一个直接证明。对每个 $\theta\in\operatorname{Hom}(E,\underline{\operatorname{Hom}}(F,G))$，如下对应一个 $\operatorname{Hom}(E\times F,G)$ 中的元素 $\phi_\theta$。对于每个 $S\in\operatorname{Ob}(\mathcal C)$，有一个关于 $S$ \Term{functorial}{函子性}的映射
\[
\theta(S):E(S)\longrightarrow\operatorname{Hom}(F\times S,G).
\]
若 $(e,f)\in E(S)\times F(S)$，则 $f$ 是\Term{morphism}{态射} $S\to F$，所以 $f\times\operatorname{id}_S$ 是\Term{morphism}{态射} $S\to F\times S$；另一方面，$\theta(S)(e)$ 是\Term{morphism}{态射} $F\times S\to G$，因此合成得到\Term{morphism}{态射}
\[
\theta(S)(e)\circ(f\times\operatorname{id}_S):S\longrightarrow G,
\]
即 $G(S)$ 的元素 $\phi_\theta(S)(e,f)$。容易核知，对应 $S\mapsto\phi_\theta(S)$ 关于 $S$ 是\Term{functorial}{函子性}的，因而定义\Term{morphism}{态射} $\phi_\theta:E\times F\to G$。映射 $\theta\mapsto\phi_\theta$ 与 $\phi\mapsto\theta_\phi$ 互为逆\Term{bijection}{双射}，留给读者核实。
\SourcePage{15}
现证明 (b)。若 $S\in\operatorname{Ob}(\mathcal C)$，由 1.7 (ii) 及将 (a) 应用于 $\mathcal C_{/S}$，有
\[
\begin{aligned}
\underline{\operatorname{Hom}}(E,\underline{\operatorname{Hom}}(F,G))(S)
&\simeq\operatorname{Hom}_S(E_S,\underline{\operatorname{Hom}}_S(F_S,G_S))\\
&\simeq\operatorname{Hom}_S(E_S\times_SF_S,G_S)\\
&\simeq\operatorname{Hom}(E\times F\times S,G)\\
&\simeq\underline{\operatorname{Hom}}(E\times F,G)(S),
\end{aligned}
\]
且这些\Term{isomorphism}{同构}关于 $S$ 是\Term{functorial}{函子性}的。
\end{proof}
\begin{SGAresult}{推论}{1.7.2}
有
\[
\operatorname{Hom}(E,\underline{\operatorname{Hom}}(F,G))\simeq\operatorname{Hom}(F,\underline{\operatorname{Hom}}(E,G)),
\]
\[
\underline{\operatorname{Hom}}(E,\underline{\operatorname{Hom}}(F,G))\simeq\underline{\operatorname{Hom}}(F,\underline{\operatorname{Hom}}(E,G)).
\]
\end{SGAresult}
特别地，令 $E=e$，并注意 $\underline{\operatorname{Hom}}(e,G)\simeq G$，有
\[
\Gamma(\underline{\operatorname{Hom}}(F,G))\simeq\operatorname{Hom}(F,G).
\]
注意，$\operatorname{Hom}$ 的合成给出\Term{functorial}{函子性}的\Term{morphism}{态射}
\[
\underline{\operatorname{Hom}}(F,G)\times\underline{\operatorname{Hom}}(G,H)\longrightarrow\underline{\operatorname{Hom}}(F,H).
\]
若 $F,G$ 是 $\widehat{\mathcal C}$ 的两个\Term{object}{对象}，用 $\operatorname{Isom}(F,G)$ 表示 $\operatorname{Hom}(F,G)$ 中由 $F$ 到 $G$ 的\Term{isomorphism}{同构}组成的\Term{subset}{子集}。由此定义 $\underline{\operatorname{Hom}}(F,G)$ 的\Term{subobject}{子对象} $\underline{\operatorname{Isom}}(F,G)$：
\[
\underline{\operatorname{Isom}}(F,G)(S)=\operatorname{Isom}(F_S,G_S).
\]
于是有\Term{isomorphism}{同构}
\[
\Gamma(\underline{\operatorname{Isom}}(F,G))\simeq\operatorname{Isom}(F,G),\qquad
\underline{\operatorname{Isom}}(F,G)\simeq\underline{\operatorname{Isom}}(G,F).
\]
在 $F=G$ 的特殊情形，置
\[
\begin{aligned}
\underline{\operatorname{End}}(F)&=\underline{\operatorname{Hom}}(F,F),&
\operatorname{End}(F)&=\operatorname{Hom}(F,F)\simeq\Gamma(\underline{\operatorname{End}}(F)),\\
\underline{\operatorname{Aut}}(F)&=\underline{\operatorname{Isom}}(F,F),&
\operatorname{Aut}(F)&=\operatorname{Isom}(F,F)\simeq\Gamma(\underline{\operatorname{Aut}}(F)).
\end{aligned}
\]
\Term{object}{对象} $\underline{\operatorname{Hom}}$、$\underline{\operatorname{Isom}}$、$\underline{\operatorname{Aut}}$、$\underline{\operatorname{End}}$ 的形成与\Term{base-change}{基变换}交换。

\begin{SGAresult}{注记}{1.7.3}
\footnote{\textsuperscript{(16)} 编者注：为了后文引用，增补了编号 1.7.3。}
注意，也可如下构造一个\Term{isomorphism}{同构}于 $\underline{\operatorname{Isom}}(F,G)$ 的\Term{object}{对象}：有一个\Term{morphism}{态射}
\[
\underline{\operatorname{Hom}}(F,G)\times\underline{\operatorname{Hom}}(G,F)\longrightarrow\underline{\operatorname{End}}(F);
\]
互换 $F,G$，可\Term{pushout}{推出}\Term{morphism}{态射}
\[
\underline{\operatorname{Hom}}(F,G)\times\underline{\operatorname{Hom}}(G,F)\longrightarrow\underline{\operatorname{End}}(F)\times\underline{\operatorname{End}}(G).
\]
另一方面，$F$ 的\Term{identity-morphism}{恒等态射}是 $\operatorname{End}(F)$ 的元素，因此定义一个\Term{morphism}{态射} $e\to\underline{\operatorname{End}}(F)$。
\SourcePage{16}
对 $G$ 作同样处理并取积，得到\Term{morphism}{态射}
\[
e\longrightarrow\underline{\operatorname{End}}(F)\times\underline{\operatorname{End}}(G).
\]
立即可见，$e$ 与 $\underline{\operatorname{Hom}}(F,G)\times\underline{\operatorname{Hom}}(G,F)$ 在 $\underline{\operatorname{End}}(F)\times\underline{\operatorname{End}}(G)$ 上的\Term{fiber-product}{纤维积}\Term{isomorphism}{同构}于 $\underline{\operatorname{Isom}}(F,G)$。
\end{SGAresult}
上述所有定义特别适用于 $F=h_X,G=h_Y$ 的情形。若 $\underline{\operatorname{Hom}}(h_X,h_Y)$ 由 $\mathcal C$ 的某个\Term{object}{对象}代表，就把这个\Term{object}{对象}记为 $\underline{\operatorname{Hom}}(X,Y)$。它具有下列性质：若 $Z\times X$ 存在，则
\[
\operatorname{Hom}(Z,\underline{\operatorname{Hom}}(X,Y))\simeq\operatorname{Hom}(Z\times X,Y).
\]
当 $\mathcal C$ 中存在积时，这条性质刻画了该\Term{object}{对象}。

同样可定义（只要它们存在）\Term{object}{对象}
\[
\underline{\operatorname{Isom}}(X,Y),\quad\underline{\operatorname{End}}(X),\quad\underline{\operatorname{Aut}}(X).
\]
这里只指出，由前面的构造，只要 $\mathcal C$ 中\Term{fiber-product}{纤维积}存在，并且 $\underline{\operatorname{Hom}}(X,Y)$、$\underline{\operatorname{Hom}}(Y,X)$、$\underline{\operatorname{End}}(X)$ 和 $\underline{\operatorname{End}}(Y)$ 存在，$\underline{\operatorname{Isom}}(X,Y)$ 就存在。

上述全部内容也适用于形如 $\mathcal C_{/S}$ 的\Term{category}{范畴}。我们将尽量用适当的\Term{signature}{符号}明确表示相应\Term{object}{对象}：例如，若 $T,T'$ 是 $\mathcal C$ 中位于 $S$ 之上的两个\Term{object}{对象}，就用 $\underline{\operatorname{Hom}}_S(T,T')$ \Term{representing-object}{表示对象} $\underline{\operatorname{Hom}}_{\mathcal C_{/S}}(T/S,T'/S)$。

\textbf{1.8. \Term{constant-object}{常对象}。} 设 $\mathcal C$ 是具有\Term{direct-sum}{直和}及\Term{fiber-product}{纤维积}、并且\Term{direct-sum}{直和}与\Term{base-change}{基变换}交换的\Term{category}{范畴}（例如\Term{scheme}{概形}\Term{category}{范畴}\footnote{\textsuperscript{(17)} 编者注：已将旧术语“\Term{prescheme}{预概形}／\Term{scheme}{概形}”替换为现行术语“\Term{scheme}{概形}／\Term{separated-scheme}{分离概形}”。}）。对于每个\Term{set}{集合} $E$ 和 $\mathcal C$ 的每个\Term{object}{对象} $S$，置
\[
E_S=\text{一个对象族 $(S_i)_{i\in E}$ 的直和，其中每个 $S_i$ 都同构于 $S$。}\tag{1.8.1}
\]
这个\Term{object}{对象}由如下公式刻画：对于每个 $T\in\operatorname{Ob}(\mathcal C)$，
\[
\operatorname{Hom}_{\mathcal C}(E_S,T)=\operatorname{Hom}(E,\operatorname{Hom}_{\mathcal C}(S,T)),\tag{1.8.2}
\]
其中右端第一个 $\operatorname{Hom}$ 取于\Term{set}{集合}\Term{category}{范畴}。

\Term{object}{对象} $E_S$ 带有到 $S$ 的\Term{canonical-projection}{典范投影}，所以 $E\mapsto E_S$ 实际上是从 $(\mathrm{Ens})$ 到 $\mathcal C_{/S}$ 的\Term{functor}{函子}。

\footnote{\textsuperscript{(18)} 编者注：在以下三段中，调整了语句顺序，并对下面假设 $(*)$ 的\Term{action}{作用}补充了一些说明。}
若 $S'\to S$ 是 $\mathcal C$ 中的\Term{arrow}{箭头}，则因\Term{direct-sum}{直和}与\Term{base-change}{基变换}交换，有
\[
E_{S'}=(E_S)_{S'}.
\]
特别地，若 $\mathcal C$ 有\Term{final-object}{终对象} $e$，则
\[
E_S=(E_e)_S.
\]
\SourcePage{17}
从 $(\mathrm{Ens})$ 到 $\mathcal C_{/S}$ 的\Term{functor}{函子} $E\mapsto E_S/S$ 与\Term{finite-product}{有限积}交换。为此只需证明
\[
E_S\times_SF_S=(E\times F)_S.\tag{$\times$}
\]
将 1.7 的结果用于 $\mathcal C_{/S}$，对于每个 $T\in\operatorname{Ob}(\mathcal C_{/S})$ 有自然\Term{isomorphism}{同构}（未注明的 $\operatorname{Hom}$ 均取于 $\mathcal C_{/S}$）：
\[
\begin{aligned}
\operatorname{Hom}((E\times F)_S,T)
&\simeq\operatorname{Hom}_{(\mathrm{Ens})}(E\times F,\operatorname{Hom}(S,T))\\
&\simeq\operatorname{Hom}_{(\mathrm{Ens})}(E,\operatorname{Hom}_{(\mathrm{Ens})}(F,\operatorname{Hom}(S,T)))\\
&\simeq\operatorname{Hom}_{(\mathrm{Ens})}(E,\operatorname{Hom}(F_S,T)),
\end{aligned}
\]
以及
\[
\begin{aligned}
\operatorname{Hom}(E_S\times_SF_S,T)
&\simeq\operatorname{Hom}(E_S,\underline{\operatorname{Hom}}(F_S,T))\\
&\simeq\operatorname{Hom}_{(\mathrm{Ens})}(E,\operatorname{Hom}(S,\underline{\operatorname{Hom}}(F_S,T))).
\end{aligned}
\]
而 $\operatorname{Hom}(S,\underline{\operatorname{Hom}}(F_S,T))\simeq\operatorname{Hom}(F_S,T)$，故得 $(\times)$。

用 $\varnothing$ 表示 $\mathcal C$ 的\Term{initial-object}{始对象}，假定图
\[
\begin{tikzcd}
\varnothing\arrow[r]\arrow[d]&S\arrow[d]\\
S\arrow[r]&S\amalg S
\end{tikzcd}\tag{$*$}
\]
是\Term{cartesian}{笛卡尔}图。（\Term{scheme}{概形}\Term{category}{范畴}就满足这一条件。）\footnote{\textsuperscript{(19)} 编者注：若 $\mathcal C$ 是只有\Term{arrow}{箭头} $A\to B$ 及 $\operatorname{id}_A,\operatorname{id}_B$ 的\Term{category}{范畴}，则 $(*)$ 不成立；此时 $B\amalg B=B$，而 $B\times_BB=B\not\simeq A$。} 那么\Term{functor}{函子} $E\mapsto E_S$ 与\Term{finite-inverse-limit}{有限投射极限}交换。

事实上，考虑到 $(\times)$，只需证明 $E\mapsto E_S$ 与\Term{fiber-product}{纤维积}交换。设 $u:E\to G$ 和 $v:F\to G$ 为两个\Term{set}{集合}映射。由于 $\mathcal C$ 中\Term{direct-sum}{直和}与\Term{base-change}{基变换}交换，有
\[
F_S\times_{G_S}E_S\simeq\coprod_{f\in F}(S_f\times_{G_S}E_S)
\simeq\coprod_{\substack{f\in F\\x\in E}}(S_f\times_{G_S}S_x).\tag{1}
\]
若 $v(f)\ne u(x)$，则在 $\mathcal C_{/S}$ 中存在一个\Term{morphism}{态射}
\[
S_f\times_{G_S}S_x\longrightarrow S_f\times_{S_{v(f)}\amalg S_{u(x)}}S_x;
\]
由假设 $(*)$，右端是 $\varnothing$。因此，
\[
S_f\times_{G_S}S_x\simeq\varnothing\quad\text{若 }v(f)\ne u(x).\tag{2}
\]
另一方面，若 $v(f)=u(x)$，则在 $\mathcal C_{/S}$ 中有\Term{morphism}{态射}
\[
S\longrightarrow S_f\times_{G_S}S_x;
\]
由于 $S\xrightarrow{\operatorname{id}}S$ 是 $\mathcal C_{/S}$ 的\Term{final-object}{终对象}，得到
\[
S_f\times_{G_S}S_x\simeq S\quad\text{若 }v(f)=u(x).\tag{3}
\]
\SourcePage{18}
合并 (1)、(2)、(3)，得到关于 $S$ \Term{functorial}{函子性}的\Term{isomorphism}{同构}
\[
F_S\times_{G_S}E_S\simeq\coprod_{F\times_GE}S=(F\times_GE)_S.
\]
形如 $E_S$ 的\Term{object}{对象}称为\Term{constant-object}{常对象}。注意，有关于 $E$ \Term{functorial}{函子性}的\Term{morphism}{态射}
\[
E\longrightarrow\Gamma(E_S/S),
\]
它将每个 $i\in E$ 映到由 $S$ 到 $S_i$ 的\Term{isomorphism}{同构}定义的 $E_S$ 在 $S$ 上的\Term{section}{截面}。假定 $\mathcal C$ 的每个\Term{object}{对象} $S$ 都满足条件 $(*)$；那么对于每个 $S\not\simeq\varnothing$，\Term{morphism}{态射} $E\to\Gamma(E_S/S)$ 是\Term{monomorphism}{单态射}。

若 $\mathcal C$ 是\Term{scheme}{概形}\Term{category}{范畴}，则 $\Gamma(E_S/S)$ 等同于从\Term{topological-space}{拓扑空间} $S$ 到\Term{set}{集合} $E$ 的\Term{locally-constant-map}{局部常值映射}\Term{set}{集合}，上述映射把 $E$ 的每个元素对应到相应的常值映射。由刚才的论述还可看出，$E_S$ 也可定义为如下\Term{functor}{函子}的代表\Term{object}{对象}：该\Term{functor}{函子}将位于 $S$ 之上的每个 $S'$ 映到从\Term{topological-space}{拓扑空间} $S'$ 到\Term{set}{集合} $E$ 的\Term{locally-constant}{局部常值}函数的\Term{set}{集合}\footnote{\textsuperscript{(20)} 编者注：还应指出，\Term{diagonal-morphism}{对角态射} $E_S\to E_S\times_SE_S=(E\times E)_S$ 是\Term{closed-immersion}{闭浸入}，即 $E_S$ 在 $S$ 上\Term{separated}{分离}。}。

\SGAsection{2}{代数结构}
给定\Term{set}{集合}\Term{category}{范畴}中的一种\Term{algebraic-structure}{代数结构}，我们要将它扩展到\Term{category}{范畴} $\mathcal C$。先讨论一个例子：群的情形。

\textbf{2.1. \Term{group-structure}{群结构}。} 沿用上一节的记号。
\begin{SGAresult}{定义}{2.1.1}
设 $G\in\operatorname{Ob}(\widehat{\mathcal C})$。在 $G$ 上给定一个 $\widehat{\mathcal C}$-\Term{group-structure}{群结构}，是指对于每个 $S\in\operatorname{Ob}(\mathcal C)$，在\Term{set}{集合} $G(S)$ 上给定一个\Term{group-structure}{群结构}，使得对于 $\mathcal C$ 中每个\Term{arrow}{箭头} $f:S'\to S''$，映射 $G(f):G(S'')\to G(S')$ 都是\Term{group-homomorphism}{群同态}。若 $G,H$ 是两个 $\widehat{\mathcal C}$-群，则从 $G$ 到 $H$ 的 $\widehat{\mathcal C}$-\Term{group-morphism}{群态射}，是一个\Term{morphism}{态射} $u\in\operatorname{Hom}(G,H)$，使得每个 $S\in\operatorname{Ob}(\mathcal C)$ 对应的\Term{set}{集合}映射 $u(S):G(S)\to H(S)$ 都是\Term{group-homomorphism}{群同态}。
\end{SGAresult}
用 $\operatorname{Hom}_{\widehat{\mathcal C}\text{-}\mathrm{Gr.}}(G,H)$ 表示从 $G$ 到 $H$ 的 $\widehat{\mathcal C}$-\Term{group-morphism}{群态射}的\Term{set}{集合}，用 $(\widehat{\mathcal C}\text{-}\mathrm{Gr.})$ 表示 $\widehat{\mathcal C}$-群\Term{category}{范畴}。

\textbf{例。} 设 $E\in\operatorname{Ob}(\widehat{\mathcal C})$；\Term{object}{对象} $\underline{\operatorname{Aut}}(E)$ 自然带有一个 $\widehat{\mathcal C}$-\Term{group-structure}{群结构}。\Term{final-object}{终对象} $e$ 带有唯一的 $\widehat{\mathcal C}$-\Term{group-structure}{群结构}，使它成为 $(\widehat{\mathcal C}\text{-}\mathrm{Gr.})$ 的\Term{final-object}{终对象}。

对于每个 $S\in\operatorname{Ob}(\mathcal C)$，令 $e_G(S)$ 为 $G(S)$ 的\Term{identity-element}{单位元}。族 $e_G(S)$ 定义元素 $e_G\in\Gamma(G)=\operatorname{Hom}(e,G)$，它是一个 $\widehat{\mathcal C}$-\Term{group-morphism}{群态射}
\SourcePage{19}
$e\to G$，称为 $G$ 的\Term{unit-section}{单位截面}。

注意，在 $G$ 上给定 $\widehat{\mathcal C}$-\Term{group-structure}{群结构}，等同于给定 $G$ 上的一个\Term{composition-law}{合成律}，即一个 $\widehat{\mathcal C}$-\Term{morphism}{态射}
\[
\pi_G:G\times G\longrightarrow G,
\]
使得对每个 $S\in\operatorname{Ob}(\mathcal C)$，$\pi_G(S)$ 都使 $G(S)$ 成为群。

同样，$f:G\to H$ 是 $\widehat{\mathcal C}$-\Term{group-morphism}{群态射}，当且仅当下图交换：
\[
\begin{tikzcd}
G\times G\arrow[r,"\pi_G"]\arrow[d,"{(f,f)}"']&G\arrow[d,"f"]\\
H\times H\arrow[r,"\pi_H"]&H.
\end{tikzcd}
\]
若 $H$ 是 $G$ 的\Term{subobject}{子对象}，且对于每个 $S\in\operatorname{Ob}(\mathcal C)$，$H(S)$ 是 $G(S)$ 的\Term{subgroup}{子群}，则它显然带有由 $G$ 的结构诱导的 $\widehat{\mathcal C}$-\Term{group-structure}{群结构}：这是使\Term{monomorphism}{单态射} $H\to G$ 成为 $\widehat{\mathcal C}$-\Term{group-morphism}{群态射}的唯一结构。带有此结构的 $\widehat{\mathcal C}$-群 $H$ 称为 $G$ 的子 $\widehat{\mathcal C}$-群。

若 $G,H$ 为两个 $\widehat{\mathcal C}$-群，则积 $G\times H$ 自然带有一个 $\widehat{\mathcal C}$-\Term{group-structure}{群结构}：对于每个 $S\in\operatorname{Ob}(\mathcal C)$，在 $G(S)\times H(S)$ 上取由 $G(S)$ 和 $H(S)$ 的\Term{group-structure}{群结构}构成的积\Term{group-structure}{群结构}。带有这一结构的 $\widehat{\mathcal C}$-群 $G\times H$ 称为 $G,H$ 的积 $\widehat{\mathcal C}$-群（事实上，它就是这两个\Term{object}{对象}在 $\widehat{\mathcal C}$-群\Term{category}{范畴}中的积）。

若 $G$ 是一个 $\widehat{\mathcal C}$-群，则对于每个 $S\in\operatorname{Ob}(\mathcal C)$，$G_S$ 是 $\widehat{\mathcal C_{/S}}$-群。若 $G,H$ 是两个 $\widehat{\mathcal C}$-群，则如下定义 $\widehat{\mathcal C}$ 的\Term{object}{对象} $\underline{\operatorname{Hom}}_{\widehat{\mathcal C}\text{-}\mathrm{Gr.}}(G,H)$：
\[
\underline{\operatorname{Hom}}_{\widehat{\mathcal C}\text{-}\mathrm{Gr.}}(G,H)(S)
=\operatorname{Hom}_{\widehat{\mathcal C_{/S}}\text{-}\mathrm{Gr.}}(G_S,H_S).
\]
（注意：$\underline{\operatorname{Hom}}_{\widehat{\mathcal C}\text{-}\mathrm{Gr.}}$ 一般不是 $\widehat{\mathcal C}$-群，更不是\Term{category}{范畴} $(\widehat{\mathcal C}\text{-}\mathrm{Gr.})$ 中的 $\underline{\operatorname{Hom}}$ \Term{object}{对象}。）

同样定义\Term{object}{对象}
\[
\underline{\operatorname{Isom}}_{\widehat{\mathcal C}\text{-}\mathrm{Gr.}}(G,H),\quad
\underline{\operatorname{End}}_{\widehat{\mathcal C}\text{-}\mathrm{Gr.}}(G),\quad
\underline{\operatorname{Aut}}_{\widehat{\mathcal C}\text{-}\mathrm{Gr.}}(G).
\]
\begin{SGAresult}{定义}{2.1.2}
设 $G\in\operatorname{Ob}(\mathcal C)$。$G$ 上的 $\mathcal C$-\Term{group-structure}{群结构}，是 $h_G\in\operatorname{Ob}(\widehat{\mathcal C})$ 上的 $\widehat{\mathcal C}$-\Term{group-structure}{群结构}。从 $\mathcal C$-群 $G$ 到 $\mathcal C$-群 $H$ 的\Term{morphism}{态射}，是一个元素 $u\in\operatorname{Hom}(G,H)\simeq\operatorname{Hom}(h_G,h_H)$，它定义从 $h_G$ 到 $h_H$ 的 $\widehat{\mathcal C}$-\Term{group-morphism}{群态射}。
\end{SGAresult}
用 $(\mathcal C\text{-}\mathrm{Gr.})$ 表示 $\mathcal C$-群\Term{category}{范畴}。注意，在 $(\mathrm{Cat})$ 中有一个\Term{cartesian-square}{笛卡尔方块}
\[
\begin{tikzcd}
(\mathcal C\text{-}\mathrm{Gr.})\arrow[r]\arrow[d]&(\widehat{\mathcal C}\text{-}\mathrm{Gr.})\arrow[d]\\
\mathcal C\arrow[r,"h"]&\widehat{\mathcal C}.
\end{tikzcd}
\]
\SourcePage{20}
所以，只要上述定义和构造中涉及的\Term{functor}{函子}（积、$\underline{\operatorname{Hom}}$ \Term{object}{对象}等）\Term{representable}{可表}，它们就立即搬到 $(\mathcal C\text{-}\mathrm{Gr.})$ 中。它们也适用于\Term{category}{范畴} $\mathcal C_{/S}$。在这种情形，我们用 $\underline{\operatorname{Hom}}_{S\text{-}\mathrm{Gr.}}$ 代替 $\underline{\operatorname{Hom}}_{\widehat{\mathcal C_{/S}}\text{-}\mathrm{Gr.}}$，等等。

\textbf{2.2.} 更一般地，若 $(T)$ 是由\Term{finite-inverse-limit}{有限投射极限}定义的、以 $n$ 个\Term{set}{集合}为\Term{underlying-set}{基集}的一种结构（例如，通过由\Term{cartesian-product}{笛卡尔积}构造的图的\Term{commutativity}{交换性}定义的\Term{monoid}{幺半群}、群、带\Term{action}{作用}的\Term{set}{集合}、环上的模、环上的 \Term{lie-algebra}{Lie 代数}等结构），则以上构造使我们能定义“$\widehat{\mathcal C}$ 的 $n$ 个\Term{object}{对象} $F_1,\ldots,F_n$ 上的 $(T)$ 型结构”这一概念：这样的结构，是对于 $\mathcal C$ 中每个 $S$，在\Term{set}{集合} $F_1(S),\ldots,F_n(S)$ 上给定一个 $(T)$ 型结构，使得对于 $\mathcal C$ 中每个\Term{arrow}{箭头} $S'\to S''$，映射族 $(F_i(S''))\to(F_i(S'))$ 都是 $(T)$ 型结构的多重\Term{homomorphism}{同态}。类似地定义 $(T)$ 型结构的\Term{morphism}{态射}，从而得到\Term{category}{范畴} $(\widehat{\mathcal C}\times\widehat{\mathcal C}\times\cdots\times\widehat{\mathcal C})^{(T)}$。\Term{fully-faithful}{全忠实函子} $h\times h\times\cdots\times h$ 于是允许通过原像定义\Term{category}{范畴} $(\mathcal C\times\mathcal C\times\cdots\times\mathcal C)^{(T)}$；又因为它与\Term{inverse-limit}{投射极限}交换，可以将 $\widehat{\mathcal C}$ 中引入的全部\Term{functorial}{函子性}性质、概念和记号搬到这一\Term{category}{范畴}中。

现在假设 $\mathcal C$ 中\Term{fiber-product}{纤维积}存在，且 $(T)$ 是一种\Term{algebraic-structure}{代数结构}，由某些\Term{cartesian-product}{笛卡尔积}之间的\Term{morphism}{态射}给定，其公理是由这些\Term{arrow}{箭头}构成的某些图的\Term{commutativity}{交换性}。因此，在 $\mathcal C$ 的一个\Term{object}{对象}族上的 $(T)$ 型结构，就由\Term{cartesian-product}{笛卡尔积}之间满足某些交换条件的若干\Term{morphism}{态射}定义。由此可见，若 $\mathcal C,\mathcal C'$ 为两个有积的\Term{category}{范畴}，且 $f:\mathcal C\to\mathcal C'$ 为与积交换的\Term{functor}{函子}，则 $\mathcal C$ 中每个带有 $(T)$ 型结构的\Term{object}{对象}族 $(F_i)$，都会使 $\mathcal C'$ 中的\Term{object}{对象}族 $(f(F_i))$ 自然带有一个 $(T)$ 型结构。每个 $\mathcal C$-群变为一个 $\mathcal C'$-群，每一对（$\mathcal C$-环，此 $\mathcal C$-环上的 $\mathcal C$-模）变为 $\mathcal C'$ 中的相应二元组，等等。

特别地，设 $\mathcal C$ 满足 1.8 的条件\footnote{\textsuperscript{(21)} 编者注：包括条件 $(*)$。}；该处定义的\Term{functor}{函子} $E\mapsto E_S$ 与\Term{finite-inverse-limit}{有限投射极限}交换，所以它把群变成 $S$-群（即 $\mathcal C_{/S}$-群）、把环变成 $S$-环，等等。

\textbf{注记。} 应指出，将以上构造用于\Term{category}{范畴} $\widehat{\mathcal C}$，确实给出我们已经定义过的概念；换言之，在 $\widehat{\mathcal C}$ 的一个\Term{object}{对象}上，当它被看作 $\mathcal C$ 上的\Term{functor}{函子}时给定 $(T)$ 型结构，等同于在此\Term{object}{对象}所定义的 $\widehat{\mathcal C}$ 上的\Term{representable-functor}{可表函子}上给定 $(T)$ 型结构\footnote{\textsuperscript{(22)} 编者注：例如，对\Term{group-structure}{群结构}，设 $G\in\operatorname{Ob}(\widehat{\mathcal C})$；若\Term{functor}{函子} $\widehat{\mathcal C}\to(\mathrm{Ens})$，$F\mapsto\operatorname{Hom}_{\widehat{\mathcal C}}(F,G)$ 带有一个\Term{group-structure}{群结构}，则其在 $\mathcal C$ 上的\Term{restriction}{限制} $X\mapsto G(X)$ 也是如此。反过来，若 $G$ 是 $\widehat{\mathcal C}$-群，则“\Term{multiplication}{乘法}”\Term{morphism}{态射} $\pi_G:G\times G\to G$ 对每个 $F\in\operatorname{Ob}(\widehat{\mathcal C})$ 都在 $\operatorname{Hom}_{\widehat{\mathcal C}}(F,G)$ 上诱导一个关于 $F$ \Term{functorial}{函子性}的\Term{group-structure}{群结构}。}。
\SourcePage{21}
我们还要讨论上述构造的两个特例：带\Term{operator-group}{算子群}的结构，以及模。

\textbf{2.3. 带\Term{operator-group}{算子群}的结构。}
\begin{SGAresult}{定义}{2.3.1}
设 $E\in\operatorname{Ob}(\widehat{\mathcal C})$ 且 $G\in\operatorname{Ob}(\widehat{\mathcal C}\text{-}\mathrm{Gr.})$。在 $E$ 上给定一个以 $\widehat{\mathcal C}$-群 $G$ 为\Term{operator-group}{算子群}的\Term{object}{对象}结构（或 $G$-\Term{object}{对象}结构），是对于每个 $S\in\operatorname{Ob}(\mathcal C)$，在 $E(S)$ 上给定一个以 $G(S)$ 为\Term{operator-group}{算子群}的\Term{set}{集合}结构，使得对于 $\mathcal C$ 中每个\Term{arrow}{箭头} $S'\to S''$，\Term{set}{集合}映射 $E(S'')\to E(S')$ 与\Term{operator}{算子}\Term{homomorphism}{同态} $G(S'')\to G(S')$ 相容。
\end{SGAresult}
像通常一样，这等同于给定一个\Term{morphism}{态射}
\[
\mu:G\times E\longrightarrow E,
\]
它对于每个 $S$ 都使 $E(S)$ 成为一个 $G(S)$-\Term{action}{作用}\Term{set}{集合}。但 $\operatorname{Hom}(G\times E,E)\simeq\operatorname{Hom}(G,\underline{\operatorname{End}}(E))$，所以 $\mu$ 定义\Term{morphism}{态射} $G\to\underline{\operatorname{End}}(E)$；立即可见，此\Term{morphism}{态射}将 $G$ 映入 $\underline{\operatorname{Aut}}(E)$，并且是 $\widehat{\mathcal C}$-\Term{group-morphism}{群态射}。因此：在 $E$ 上给定一个以 $\widehat{\mathcal C}$-群 $G$ 为\Term{operator-group}{算子群}的\Term{object}{对象}结构，等同于给定一个 $\widehat{\mathcal C}$-\Term{group-morphism}{群态射}
\[
\rho:G\longrightarrow\underline{\operatorname{Aut}}(E).
\]
特别地，每个元素 $g\in G(S)$ 都定义\Term{functor}{函子} $E_S$ 的\Term{automorphism}{自同构} $\rho(g)$，即 $E\times h_S$ 的一个与\Term{projection}{投影} $E\times h_S\to h_S$ 交换的\Term{automorphism}{自同构}；尤其对于每个 $S'\to S$，它给出\Term{set}{集合} $E(S')$ 的一个\Term{automorphism}{自同构}。

\begin{SGAresult}{定义}{2.3.2}
用 $E^G$ 表示如下定义的 $E$ 的\Term{subobject}{子对象}：
\[
E^G(S)=\{x\in E(S)\mid \text{对每个 }S'\to S,\ x_{S'}\text{ 在 }G(S')\text{ 下不变}\},
\]
其中 $x_{S'}$ 表示 $x$ 在 $E(S)\to E(S')$ 下的像。
\end{SGAresult}
于是，$E^G$（“$G$ 的\Term{invariants}{不变量}\Term{subobject}{子对象}”）是 $E$ 中 $G$ 在其上\Term{trivial-action}{平凡作用}的最大\Term{subobject}{子对象}。

\begin{SGAresult}{定义}{2.3.3}
设 $F$ 是 $E$ 的\Term{subobject}{子对象}。用 $\operatorname{Norm}_GF$ 和 $\operatorname{Centr}_GF$ 表示 $G$ 的如下两个子 $\widehat{\mathcal C}$-群：
\[
\begin{aligned}
(\operatorname{Norm}_GF)(S)
&=\{g\in G(S)\mid\rho(g)F_S=F_S\}\\
&=\{g\in G(S)\mid\rho(g)F(S')=F(S')\text{，对每个 }S'\to S\},\\
(\operatorname{Centr}_GF)(S)
&=\{g\in G(S)\mid\rho(g)|_{F_S}=\operatorname{id}\}\\
&=\{g\in G(S)\mid\rho(g)|_{F(S')}=\operatorname{id}\text{，对每个 }S'\to S\}.
\end{aligned}
\]
\footnote{\textsuperscript{(23)} 编者注：已订正原文，将包含关系 $\rho(g)F_S\subset F_S$ 改成等式，以确保 $\operatorname{Norm}_G(F)$ 确为群（关于“\Term{transporter}{输运子}”与“\Term{strict-transporter}{严格输运子}”一致的条件，参见 VI\textsubscript{B} 6.4）。}
这里，$\rho(g)$ 后面的竖线表示\Term{restriction}{限制}。
\end{SGAresult}
\SourcePage{22}
\begin{SGAresult}{补说}{2.3.3.1}
\footnote{\textsuperscript{(24)} 编者注：增补了补说 2.3.3.1 及注记 2.3.3.2。}
特别地，设 $x\in\Gamma(E)$，即（参见 1.2）一个元素族 $x_S\in E(S)$，$S\in\operatorname{Ob}(\mathcal C)$，使得对每个\Term{arrow}{箭头} $f:S'\to S$，有 $E(f)(x_S)=x_{S'}$（若 $\mathcal C$ 有\Term{final-object}{终对象} $S_0$，则 $\Gamma(E)=E(S_0)$）。于是 $x$ 定义 $E$ 的一个\Term{subfunctor}{子函子}，记为 $x$，并且 $\operatorname{Norm}_Gx=\operatorname{Centr}_Gx$。此\Term{functor}{函子}记作 $\operatorname{Stab}_G(x)$，称为 $x$ 的\Term{stabilizer}{稳定子}；故对于每个 $S\in\operatorname{Ob}(\mathcal C)$，有
\[
\operatorname{Stab}_G(x)(S)=\{g\in G(S)\mid\rho(g)x_S=x_S\}.
\]
假定 $\mathcal C$ 中\Term{fiber-product}{纤维积}存在；若 $G=h_G$（分别地，$E=h_E$），其中 $G$ 是 $\mathcal C$-群（分别地，$E\in\operatorname{Ob}(\mathcal C)$），且 $\mathcal C$ 有\Term{final-object}{终对象} $S_0$，因而 $x$ 是一个\Term{morphism}{态射} $S_0\to E$，则 $\operatorname{Stab}_G(x)$ 由\Term{fiber-product}{纤维积} $G\times_ES_0$ 代表；这里 $G\to E$ 是 $\operatorname{id}_G\times x:G=G\times S_0\to G\times E$ 与 $\mu:G\times E\to E$ 的合成。
\end{SGAresult}
\begin{SGAresult}{注记}{2.3.3.2}
\textsuperscript{(24)} $E^G$、$\operatorname{Norm}_GF$ 和 $\operatorname{Centr}_GF$ 的形成与\Term{base-change}{基变换}交换，即对每个 $S\in\operatorname{Ob}(\mathcal C)$ 有
\[
(E^G)_S=(E_S)^{G_S},\qquad
(\operatorname{Norm}_GF)_S\simeq\operatorname{Norm}_{G_S}F_S,\qquad
(\operatorname{Centr}_GF)_S\simeq\operatorname{Centr}_{G_S}F_S.
\]
\end{SGAresult}
\begin{SGAresult}{定义}{2.3.4}
若 $G$ 是 $\mathcal C$-群，$\mathbf E$ 是 $\widehat{\mathcal C}$ 的\Term{object}{对象}（分别地，$E$ 是 $\mathcal C$ 的\Term{object}{对象}），则 $\mathbf E$（分别地，$E$）上的 $G$-\Term{object}{对象}结构，是 $\mathbf E$（分别地，$h_E$）上的 $h_G$-\Term{object}{对象}结构。
\end{SGAresult}
由此定义，只要上述概念和记号仅涉及\Term{representable-functor}{可表函子}，它们就全部搬到 $\mathcal C$ 中：例如，若 $\operatorname{Norm}_{h_G}(h_F)$ \Term{representable}{可表}，则存在且仅存在一个 $G$ 的\Term{subobject}{子对象}代表它；该\Term{subobject}{子对象}于是是 $G$ 的子 $\mathcal C$-群，记为 $\operatorname{Norm}_G(F)$，等等。

\begin{SGAresult}{定义}{2.3.5}
a) 若 $\widehat{\mathcal C}$-群 $H$ 带有一个 $G$-\Term{object}{对象}结构，使得对每个 $g\in G(S)$，由 $g$ 定义的 $H(S)$ 的\Term{automorphism}{自同构}都是群\Term{automorphism}{自同构}，就称 $\widehat{\mathcal C}$-群 $G$ \Term{action}{作用}于 $\widehat{\mathcal C}$-群 $H$。

这等同于说，对每个 $g\in G(S)$，$H_S$ 的\Term{automorphism}{自同构} $\rho(g)$ 是 $\widehat{\mathcal C_{/S}}$-群\Term{automorphism}{自同构}；也等同于说，$\widehat{\mathcal C}$-\Term{group-morphism}{群态射} $G\to\underline{\operatorname{Aut}}(H)$ 将 $G$ 映入 $\underline{\operatorname{Aut}}_{\widehat{\mathcal C}\text{-}\mathrm{Gr.}}(H)$。

b) 在上述情形，积 $H\times G$ 上存在唯一的 $\widehat{\mathcal C}$-\Term{group-structure}{群结构}，使得对于每个 $S$，$(H\times G)(S)$ 是相对于给定的 $G(S)$ 在 $H(S)$ 上的\Term{action}{作用}的群 $H(S)$ 与 $G(S)$ 的\Term{semidirect-product}{半直积}。将这一 $\widehat{\mathcal C}$-群记为 $H\cdot G$，称为 $H$ 与 $G$ 的\Term{semidirect-product}{半直积}。因此，按照定义，
\[
(H\cdot G)(S)=H(S)\cdot G(S).
\]
\end{SGAresult}
设 $G$ 为 $\widehat{\mathcal C}$-群。对于 $\mathcal C$ 中每个\Term{arrow}{箭头} $S'\to S$ 以及每个 $g\in G(S)$，令 $\operatorname{Int}(g)$ 为由 $\operatorname{Int}(g)h=ghg^{-1}$ 定义的 $G(S')$ 的\Term{automorphism}{自同构}。该定义扩充为 $\widehat{\mathcal C}$-\Term{group-morphism}{群态射}
\[
\operatorname{Int}:G\longrightarrow\underline{\operatorname{Aut}}_{\widehat{\mathcal C}\text{-}\mathrm{Gr.}}(G)\subset\underline{\operatorname{Aut}}(G).
\]
\SourcePage{23}
于是定义 2.3.3 适用，且对 $G$ 的每个\Term{subobject}{子对象} $E$，有 $G$ 的子 $\widehat{\mathcal C}$-群
\[
\operatorname{Norm}_G(E)\quad\text{和}\quad\operatorname{Centr}_G(E).
\]
\begin{SGAresult}{定义}{2.3.6}
$G$ 的子 $\widehat{\mathcal C}$-群 $\operatorname{Centr}_G(G)$ 称为 $G$ 的\Term{center}{中心}，记为 $\operatorname{Centr}(G)$。若 $\operatorname{Centr}(G)=G$，或\Term{equivalence}{等价}地，对每个 $S$，$G(S)$ 都交换，就称 $G$ 交换。

若 $\operatorname{Norm}_G(H)=G$，或\Term{equivalence}{等价}地，对每个 $S$，$H(S)$ 都是 $G(S)$ 的\Term{normal-subgroup}{正规子群}，就称 $G$ 的子 $\widehat{\mathcal C}$-群 $H$ 在 $G$ 中\Term{normal-subgroup}{正规}\footnote{\textsuperscript{(25)} 编者注：此外，若 $\operatorname{Centr}_G(H)=G$，或\Term{equivalence}{等价}地，对每个 $S$，$H(S)$ 都是 $G(S)$ 的\Term{central-subgroup}{中心子群}，就称 $H$ 在 $G$ 中\Term{center}{中心}。}。
\end{SGAresult}
\begin{SGAresult}{定义}{2.3.6.1}
\footnote{\textsuperscript{(26)} 编者注：增补了定义 2.3.6.1 和 2.3.6.2。}
设 $f:G\to G'$ 为 $\widehat{\mathcal C}$-\Term{group-morphism}{群态射}。由下式对每个 $S\in\operatorname{Ob}(\mathcal C)$ 定义的 $G$ 的子 $\widehat{\mathcal C}$-群，称为 $f$ 的\Term{kernel}{核}，记为 $\operatorname{Ker}f$：
\[
(\operatorname{Ker}f)(S)=\{x\in G(S)\mid f(S)(x)=1\}=\operatorname{Ker}f(S).
\]
它是一个\Term{normal}{正规}子 $\widehat{\mathcal C}$-群。
\end{SGAresult}
注意，若 $G=h_G$ 且 $G'=h_{G'}$，$\mathcal C$ 有\Term{final-object}{终对象} $S_0$，且 $\mathcal C$ 中\Term{fiber-product}{纤维积}存在，则 $\operatorname{Ker}(f)$ 由 $S_0\times_{G'}G$ 代表。

\begin{SGAresult}{定义}{2.3.6.2}
\textsuperscript{(26)} 设 $E\in\widehat{\mathcal C}$，且 $G$ 是\Term{action}{作用}于 $E$ 的 $\widehat{\mathcal C}$-群。若\Term{morphism}{态射} $G\to\underline{\operatorname{Aut}}(E)$ 的核是平凡的，就称 $G$ 在 $E$ 上的\Term{action}{作用}是\Term{faithful-action}{忠实}的；即对于每个 $S\in\operatorname{Ob}(\mathcal C)$ 和 $g\in G(S)$，条件“对每个 $S'\to S$ 和 $x\in E(S')$ 均有 $g_{S'}\cdot x=x$”蕴含 $g=1$。
\end{SGAresult}
初等群论的许多定义和命题都容易移植到这里。只列出以下后文将用到的一条。
\begin{SGAresult}{命题}{2.3.7}
设 $f:W\to G$ 为 $\widehat{\mathcal C}$-\Term{group-morphism}{群态射}。置 $H(S)=\operatorname{Ker}f(S)$。设 $u:G\to W$ 为 $f$ 的一个\Term{section}{截面}，且为 $\widehat{\mathcal C}$-\Term{group-morphism}{群态射}（因而必为\Term{monomorphism}{单态射}）。那么 $W$ 可等同于 $H$ 与 $G$ 的\Term{semidirect-product}{半直积}，其中 $G$ 在 $H$ 上的\Term{action}{作用}由
\[
(g,h)\longmapsto\operatorname{Int}(u(g))h
\]
定义，这里 $g\in G(S)$、$h\in H(S)$、$S\in\operatorname{Ob}(\mathcal C)$。
\end{SGAresult}
这些定义和命题像通常一样全部搬到 $\mathcal C$ 中。特别地，当\Term{cartesian-product}{笛卡尔积} $H\times G$ 存在时，可定义两个 $\mathcal C$-群 $H,G$（$G$ \Term{action}{作用}于 $H$）的\Term{semidirect-product}{半直积}，并且命题 2.3.7 有如下对应\Term{form}{形式}。
\begin{SGAresult}{命题}{2.3.8}
设 $H\xrightarrow{i}W\xrightarrow{f}G$ 是 $\mathcal C$-\Term{group-morphism}{群态射}列，使得对每个 $S\in\operatorname{Ob}(\mathcal C)$，$(H(S),i(S))$ 是 $f(S):W(S)\to G(S)$ 的一个核。设 $u:G\to W$ 是 $f$ 的一个\Term{section}{截面}，且为 $\mathcal C$-\Term{group-morphism}{群态射}。那么 $W$ 可等同于 $H$ 与 $G$ 的\Term{semidirect-product}{半直积}，其中 $G$ 在 $H$ 上的\Term{action}{作用}满足：若 $S\in\operatorname{Ob}(\mathcal C)$、$g\in G(S)$、$h\in H(S)$，则
\[
\operatorname{Int}(u(g))i(h)=i({}^gh).
\]
\end{SGAresult}
\SourcePage{24}
\SGAsection{3}{$O$-模范畴与 $G$-$O$-模范畴}
\begin{SGAresult}{定义}{3.1}
设 $O,F$ 是 $\widehat{\mathcal C}$ 的两个\Term{object}{对象}。若对于每个 $S\in\operatorname{Ob}(\mathcal C)$，在 $O(S)$ 上给定环结构，并在 $F(S)$ 上给定该环上的模结构，使得对于 $\mathcal C$ 中的每个\Term{arrow}{箭头} $S'\to S''$，$O(S'')\to O(S')$ 是\Term{ring-homomorphism}{环同态}，而 $F(S'')\to F(S')$ 是与该\Term{ring-homomorphism}{环同态}相容的\Term{abelian-group}{阿贝尔群}\Term{homomorphism}{同态}，就称 $F$ 是 $\widehat{\mathcal C}$-环 $O$ 上的 $\widehat{\mathcal C}$-模，简称 $O$-模。
\end{SGAresult}
固定 $O$ 后，按通常方式定义 $O$-模 $F,F'$ 的\Term{morphism}{态射}，从而得到\Term{abelian-group}{交换群} $\operatorname{Hom}_O(F,F')$，以及记作 $(O\text{-}\mathrm{Mod.})$ 的 $O$-\Term{module-category}{模范畴}。

\begin{SGAresult}{\Term{lemma}{引理}}{3.1.1}
\footnote{\textsuperscript{(27)} 编者注：增补了陈述 3.1.1 和 3.1.2，以说明\Term{category}{范畴} $(O\text{-}\mathrm{Mod.})$ 是\Term{abelian-category}{阿贝尔范畴}且满足公理 (AB 5)，并且当\Term{category}{范畴} $\mathcal C$ 为小时具有\Term{enough-injectives}{足够多的内射对象}。}
$(O\text{-}\mathrm{Mod.})$ 带有“逐自变量”定义的\Term{abelian-category}{阿贝尔范畴}结构。此外，$(O\text{-}\mathrm{Mod.})$ 满足公理 (AB 5)（参见 [Gr57，1.5]），即任意\Term{direct-sum}{直和}存在，并且，若 $M$ 是 $O$-模，$N$ 是其\Term{submodule}{子模}，$(F_i)_{i\in I}$ 是 $M$ 的一个\Term{filtered}{滤过}递增\Term{submodule}{子模}族，则
\[
\bigcup_{i\in I}(F_i\cap N)=\left(\bigcup_{i\in I}F_i\right)\cap N.
\]
\end{SGAresult}
\begin{proof}
设 $f:F\to F'$ 为 $O$-\Term{module-morphism}{模态射}。令每个 $S\in\operatorname{Ob}(\mathcal C)$ 对应于 $(\operatorname{Ker}f)(S)=\operatorname{Ker}f(S)$（分别地，$\cdots$），以此定义 $O$-模 $\operatorname{Ker}f$（分别地，$\operatorname{Im}f$ 和 $\operatorname{Coker}f$）。那么 $\operatorname{Ker}f$（分别地，$\operatorname{Coker}f$）是 $f$ 的一个核（分别地，\Term{cokernel}{余核}），并且有 $O$-模\Term{isomorphism}{同构} $F/\operatorname{Ker}f\xrightarrow{\sim}\operatorname{Im}f$。这证明 $(O\text{-}\mathrm{Mod.})$ 是\Term{abelian-category}{阿贝尔范畴}。

任意\Term{direct-sum}{直和}存在，并且“逐自变量”定义。最后，若 $M$ 是 $O$-模，$N$ 是其\Term{submodule}{子模}，$(F_i)_{i\in I}$ 是 $M$ 的一个\Term{filtered}{滤过}递增\Term{submodule}{子模}族，则包含关系
\[
\bigcup_{i\in I}(F_i\cap N)\subset\left(\bigcup_{i\in I}F_i\right)\cap N
\]
是等式：事实上，若 $S\in\operatorname{Ob}(\mathcal C)$ 且 $x\in N(S)\cap\bigcup_iF_i(S)$，则存在 $i\in I$，使得 $x\in N(S)\cap F_i(S)$。
\end{proof}
\begin{SGAresult}{命题}{3.1.2}
\textsuperscript{(27)} 假定\Term{category}{范畴} $\mathcal C$ 为小，即 $\operatorname{Ob}(\mathcal C)$ 是\Term{set}{集合}。那么 $O$ 是 $(O\text{-}\mathrm{Mod.})$ 的一个\Term{generator}{生成元}；因此，$(O\text{-}\mathrm{Mod.})$ 具有足够多的\Term{injective-object}{内射对象}。
\end{SGAresult}
\begin{proof}
设 $F$ 是 $O$-模。对于每个 $S\in\operatorname{Ob}(\mathcal C)$，令 $F_0(S)$ 是 $O(S)$-模 $F(S)$ 的一个\Term{generating-family}{生成元系}。由于假设 $\operatorname{Ob}(\mathcal C)$ 是\Term{set}{集合}，可以考虑\Term{set}{集合} $I=\coprod_{S\in\operatorname{Ob}(\mathcal C)}F_0(S)$；于是有\Term{epimorphism}{满态射}
\[
O^{\oplus I}\twoheadrightarrow F.
\]
\SourcePage{25}
这表明 $O$ 是 $(O\text{-}\mathrm{Mod.})$ 的一个\Term{generator}{生成元}（参见 [Gr57，1.9.1]）。由于 $(O\text{-}\mathrm{Mod.})$ 满足 (AB 5)，由 [Gr57，1.10.1] 得到 $(O\text{-}\mathrm{Mod.})$ 具有\Term{enough-injectives}{足够多的内射对象}。
\end{proof}
\begin{SGAresult}{注记}{3.1.3}
若 $O'$ 是由 $O'(S)=\mathbb Z$ 定义的 $\widehat{\mathcal C}$-环（不要将它与\Term{constant-object}{常对象} $\mathbb Z$ 关联的\Term{functor}{函子}混淆），则 $O'$-\Term{module-category}{模范畴}\Term{isomorphism}{同构}于交换 $\widehat{\mathcal C}$-群的\Term{category}{范畴}。
\end{SGAresult}
\begin{SGAresult}{定义}{3.1.4}
注意，若 $F$ 是 $O$-模，则对每个 $S\in\operatorname{Ob}(\mathcal C)$，$F_S$ 是 $O_S$-模。因此可由
\[
\underline{\operatorname{Hom}}_O(F,F')(S)=\operatorname{Hom}_{O_S}(F_S,F'_S)
\]
定义阿贝尔 $\widehat{\mathcal C}$-群 $\underline{\operatorname{Hom}}_O(F,F')$。同样定义\Term{object}{对象}
\[
\underline{\operatorname{Isom}}_O(F,F'),\quad\underline{\operatorname{End}}_O(F),\quad\underline{\operatorname{Aut}}_O(F),
\]
其中最后一个在\Term{automorphism}{自同构}的合成所诱导的结构下是 $\widehat{\mathcal C}$-群。
\end{SGAresult}
\begin{SGAresult}{定义}{3.2}
设 $O$ 为 $\widehat{\mathcal C}$-环，$F$ 为 $O$-模，$G$ 为 $\widehat{\mathcal C}$-群。$F$ 上的 $G$-$O$-模结构，是一个 $G$-\Term{object}{对象}结构，使得对于每个 $S\in\operatorname{Ob}(\mathcal C)$ 和每个 $g\in G(S)$，由 $g$ 定义的 $F(S)$ 的\Term{automorphism}{自同构}都是其 $O(S)$-模结构的\Term{automorphism}{自同构}。
\end{SGAresult}
这等同于说，2.3 中定义的 $\widehat{\mathcal C}$-\Term{group-morphism}{群态射}
\[
\rho:G\longrightarrow\underline{\operatorname{Aut}}(F)
\]
将 $G$ 映入 $\underline{\operatorname{Aut}}(F)$ 的子 $\widehat{\mathcal C}$-群 $\underline{\operatorname{Aut}}_O(F)$。因此，在 $O$-模 $F$ 上给定一个 $G$-$O$-模结构，等同于给定一个 $\widehat{\mathcal C}$-\Term{group-morphism}{群态射}
\[
\rho:G\longrightarrow\underline{\operatorname{Aut}}_O(F).
\]
自然地定义\Term{abelian-group}{阿贝尔群} $\operatorname{Hom}_{G\text{-}O}(F,F')$，从而得到 $G$-$O$-模的\Term{additive-category}{加性范畴}，记为 $(G\text{-}O\text{-}\mathrm{Mod.})$。

\begin{SGAresult}{注记}{3.2.1}
读者会注意到，$(G\text{-}O\text{-}\mathrm{Mod.})$ 也可定义为 $O[G]$-\Term{module-category}{模范畴}，其中 $O[G]$ 是 $\widehat{\mathcal C}$-群 $G$ 在 $\widehat{\mathcal C}$-环 $O$ 上的\Term{algebra}{代数}，其定义是清楚的\footnote{\textsuperscript{(28)} 编者注：增补了下句；它将在第 5 节中用到。}。因此，由 3.1.1 和 3.1.2，$(G\text{-}O\text{-}\mathrm{Mod.})$ 是满足公理 (AB 5) 的\Term{abelian-category}{阿贝尔范畴}；此外，若 $\mathcal C$ 为小，$(G\text{-}O\text{-}\mathrm{Mod.})$ 具有\Term{enough-injectives}{足够多的内射对象}。
\end{SGAresult}
本节所有构造都立即可专门用于 $G$（或 $O$，或二者）由 $\mathcal C$ 中\Term{object}{对象}代表的情形；这些\Term{object}{对象}自然带有相应的\Term{algebraic-structure}{代数结构}。

我们简略讨论了本讨论班后文遇到的主要\Term{algebraic-structure}{代数结构}。至于其他结构（例如 $O$-\Term{lie-algebra}{Lie 代数}结构），
\SourcePage{26}
我们认为，本节已经给了读者足够的例子，使他能在每个具体情形中，自己运用 2.2 所勾勒的一般机制。

现在将把刚才的工作应用于记为 $(\mathrm{Sch})$ 的\Term{scheme}{概形}\Term{category}{范畴}，更一般地，应用于\Term{category}{范畴} $(\mathrm{Sch})_{/S}$（也记为 $(\mathrm{Sch}/S)$）。
% 本块到此截止；以下原文为第 4 节“Structures algébriques dans la catégorie des schémas”。
